\documentclass[11pt,reqno]{amsart}
\usepackage[T1]{fontenc}
\usepackage[utf8]{inputenc}
\usepackage{lmodern,amsmath,amssymb,amsthm,mathtools}
\usepackage[margin=1.05in]{geometry}
\usepackage{microtype,enumitem,xcolor}
\usepackage[hyphens]{url}
\usepackage[unicode,colorlinks=true,linkcolor=blue!45!black,citecolor=blue!45!black,urlcolor=blue!45!black,pdfauthor={OpenAI},pdftitle={Integrable metrics and effectivity with controlled boundary}]{hyperref}
\numberwithin{equation}{section}
\newtheorem{theorem}{Theorem}[section]
\newtheorem{proposition}[theorem]{Proposition}
\newtheorem{lemma}[theorem]{Lemma}
\newtheorem{corollary}[theorem]{Corollary}
\theoremstyle{definition}
\newtheorem{definition}[theorem]{Definition}
\newtheorem{remark}[theorem]{Remark}
\newcommand{\Q}{\mathbb Q}
\newcommand{\R}{\mathbb R}
\newcommand{\C}{\mathbb C}
\DeclareMathOperator{\Supp}{Supp}
\DeclareMathOperator{\ord}{ord}
\DeclareMathOperator{\rk}{rk}
\DeclareMathOperator{\Pic}{Pic}
\DeclareMathOperator{\trdeg}{trdeg}
\DeclareMathOperator{\psef}{Psef}
\newcommand{\OO}{\mathcal O}
\DeclareMathOperator{\divisor}{div}
\title{Integrable metrics and effectivity with controlled boundary}
\author{OpenAI}
\date{September 26, 2026}
\begin{document}
\raggedbottom
\begin{abstract}
Let $Y$ be smooth projective and let $C$ be an effective integral divisor.
If $-K_Y+C$ admits a metric with a global strict curvature lower bound
for which the canonical section of $C$ is locally square integrable,
every pseudoeffective rational divisor becomes rationally effective
after an effective correction supported on $C$.
For a smooth projective $X$ with smoothly semipositive $L=-K_X$,
sections of $m_jL-P$ at unbounded positive exponents, with any fixed
pseudoeffective Cartier error $P$, yield a section of a positive
multiple of $L$.
\end{abstract}
\maketitle
\tableofcontents
\section{Introduction}

A pseudoeffective divisor need not have a section in any positive
multiple.  This paper studies an analytic condition that makes such a
divisor effective up to rational linear equivalence, after a correction
on a prescribed boundary.  The condition is integrability of the
canonical divisor section for a positively curved anticanonical metric.
Its usefulness is birational: a correction on the base of a fibration
can disappear after pullback and pushforward to the original variety.

If the squared norm in a local frame of a singular metric $h$ is
$e^{-\phi}$, its multiplier ideal $\mathcal J(h)$ consists of the
holomorphic germs $u$ with $|u|^2e^{-\phi}$ locally integrable.
Thus $\mathcal J(h)\supset\mathcal O_Y(-C)$ means precisely that
the canonical section $s_C$ of the effective integral divisor $C$
has locally finite squared norm.  We call a rational divisor
\emph{rationally effective} if it is rationally linearly equivalent
to an effective rational divisor.  A rational divisor is
\emph{pseudoeffective} when its numerical class lies in the closure
of the cone generated by effective divisors.  On a smooth projective
variety this is equivalent to the existence of a singular Hermitian
metric with semipositive curvature current on its rational line bundle.

\begin{theorem}[Effectivity modulo a controlled boundary]
\label{thm:controlled-boundary}
Let $Y$ be a smooth projective complex variety of positive dimension
and let $C\ge0$ be an integral divisor.  Suppose $-K_Y+C$ has a
singular Hermitian metric $h$ whose curvature dominates a K\"ahler
form and for which
\[
 \mathcal J(h)\supset\mathcal O_Y(-C).
\]
For every pseudoeffective rational divisor $M$ on $Y$, there is an
effective rational divisor $C_0$ supported on $C$ such that
$M+C_0$ is rationally effective.
\end{theorem}

The support conclusion distinguishes the theorem from numerical
nonvanishing.  It permits the correction to be removed under a
specified birational map.  The analytic assumption also distinguishes
two notions of integrability: $s_C$ may have finite norm even when a
nonvanishing local frame does not.  The stronger, unweighted condition
will yield a separate finite-generation result.

\subsection{The boundary construction and its antecedents}

An effective rational boundary $\Delta$ is \emph{klt}
(Kawamata log terminal) if, on a log resolution $\sigma:Y'\to Y$,
every coefficient of
$\sigma^*(K_Y+\Delta)-K_{Y'}$ is less than one.
For a simple normal crossing boundary on a smooth variety, this
amounts to each boundary coefficient being less than one.

The proof connects multiplier-ideal approximation with big-boundary
nonvanishing.  Nadel vanishing~\cite{Nadel1990}, followed by very ample
regularity, makes the relevant multiplier-ideal systems globally
generated.  Point extension of Ohsawa--Takegoshi~\cite{OhsawaTakegoshi1987},
in the weighted Bergman approximation developed by
Demailly~\cite{Demailly2015}, compares the fixed divisor of a complete
system with the metric's singularities.  On a resolution
$\sigma:Y'\to Y$, weighted integrability produces a signed divisor
$\Psi=\Psi^+-\Psi^-$ whose coefficients are strictly below one.
Its negative part lies only over $C$ and the exceptional locus.

The approximation leaves an ample rational class $A$ in reserve.
For a prescribed pseudoeffective $M$, choose a small rational $t>0$
so that $A+tM$ remains ample.  General free and ample representatives
then turn the positive part of the signed divisor into an
\emph{effective big klt} boundary $\Delta$, with
\[
 K_{Y'}+\Delta\sim_{\mathbb Q}t\sigma^*M+\Psi^-.
\]
This is the pair to which the nonvanishing theorem of
Birkar--Cascini--Hacon--McKernan~\cite[Theorem D]{BCHM2010} applies.
A finite rational linear system upgrades its real-linear conclusion
to rational linear equivalence.  Pushforward gives the correction
$\sigma_*\Psi^-/t$, supported on $C$.
Sections~\ref{sec:approximation} and~\ref{sec:effectivity} give this
construction.  The latter also retains an analytic perturbation
argument: strong openness of Guan--Zhou~\cite{GuanZhou2015} and
small-exponent integrability of Skoda~\cite{Skoda1972} permit a small
pseudoeffective perturbation of the metric while preserving the
integrability of the distinguished section.

A different consequence begins with an unweighted integrable
semipositive metric that is smooth and strictly positive on one ordinary
open set.
Boucksom's volume criterion~\cite{Boucksom2002} first gives bigness;
a small mixture with a strict metric then gives a global curvature
lower bound while preserving integrability.  Section~\ref{sec:correction-ring}
makes this preparation explicit and permits perturbation in finitely
many arbitrary rational directions.  The resulting klt boundaries
fit the multigraded adjoint finite-generation results of
Cascini--Lazi\'c~\cite{CasciniLazic2012} and
Corti--Lazi\'c~\cite[Theorem 3.2(1)]{CortiLazic2013}.
More precisely, if $-K_Y+A_0$ has this unweighted metric, then for
every finite list of rational divisors $P_1,\ldots,P_s$ there are a
common rational $\epsilon>0$ and Cartier multiples
$B_i=k_0(A_0+\epsilon P_i)$ whose multigraded section ring is finitely
generated (Theorem~\ref{thm:correction-ring}).  This controls the
sections of specified Cartier divisors, not merely a numerical cone.

\subsection{Anticanonical nonvanishing and fixed-error conversion}

For a smooth projective variety $X$, anticanonical nonvanishing asks
whether some $H^0(X,-mK_X)$, $m>0$, is nonzero when $-K_X$ is nef.
Numerical effectivity is weaker: an effective divisor in the numerical
class need not supply a section of the actual line bundle.
For pairs $(X,\Delta)$, the corresponding divisor is
$-(K_X+\Delta)$, the anti-log-canonical divisor.
Lazi\'c--Matsumura--Peternell--Tsakanikas--Xie
\cite[Theorem~A]{LMPTX2023} established numerical nonvanishing for
projective log canonical threefold pairs with nef anti-log-canonical
divisor, assuming either that the variety is $\mathbb Q$-factorial
or that it has rational singularities.  M\"uller
\cite[Corollary~B]{Muller2025} proved nonvanishing for projective klt
threefold pairs with nef anti-log-canonical divisor.  His higher-dimensional
theorem assumes, in addition to nefness, that the anti-log-canonical
divisor of the induced pair on the general fiber of the maximal
rationally connected fibration is semiample \cite[Theorem~A]{Muller2025}.

Our principal conversion result uses smooth metric semipositivity and
works in every dimension.  Let $X$ be smooth connected projective,
let $L=-K_X$ carry a smooth semipositive Hermitian metric, and let $P$
be any fixed pseudoeffective Cartier divisor.  Then
Theorem~\ref{conv:fixed-error} gives the implication
\[
 H^0(X,m_jL-P)\ne0\ \text{for unbounded positive }m_j
 \quad\Longrightarrow\quad
 H^0(X,mL)\ne0\ \text{for some }m>0.
\]
The fixed error $P$ may itself have no effective representative.
A single Iitaka base converts the effective divisors into a
pseudoeffective rational divisor on that base.  The controlled-boundary
theorem makes this divisor rationally effective with correction on
base primes whose pullbacks disappear over $X$.  A rational function
from the base then corrects the original rational section without
introducing a pole on $X$.

Twisted differential forms supply the sequence through the determinant
method of Lazi\'c--Peternell \cite[Lemma~4.1]{LP2018}, adapted in
\cite[Lemma~5.1]{LMPTX2023}.  Ou's generic tangent-nefness theorem
\cite[Theorem~1.4]{Ou2023}, in the cotangent-subsheaf form
\cite[Theorem~4.1]{LMPTX2023}, makes the negative determinant
pseudoeffective.  Our extraction argument starts with positive
unbounded exponents and also includes numerically trivial $L$.
Thus nonzero sections of $(\Omega_X^1)^{\otimes p}\otimes m_jL$
for one fixed $p\ge0$ and unbounded positive $m_j$ give an
anticanonical section by Corollary~\ref{conv:tensor-conversion}.
The same holds for exterior forms.  This is an ordinary conversion;
invariant sections require further control of characters.

\subsection{Volume metrics and the distinct geometric refinements}

Let $X$ be smooth projective and let $L=-K_X$ carry a smooth
semipositive metric.  In a resolved rational fibration
\[
 X\xleftarrow{\ \pi\ } Z\xrightarrow{\ g\ }Y,
\]
assume $Z,Y$ are smooth projective, $\pi$ is birational, and $g$ is
surjective.  Write $K_{Z/X}=K_Z-\pi^*K_X$ for the relative canonical
(Jacobian) divisor.  Its canonical section belongs to $K_Z+\pi^*L$.
When
$\operatorname{rank}g_*\mathcal O_Z(K_{Z/X})=1$, its squared fiberwise
$L^2$ norm is a volume density $\rho$ on the base.  On the smooth
base-change open where this section generates the adjoint line,
$-\log\rho$ is a metric on $-K_Y$.
Berndtsson's direct-image theory supplies its positivity
\cite{Berndtsson2009,Berndtsson2011}.  Extending it over the boundary
with enough integrability is the geometric task preceding effectivity.

The main conversion proof uses a singular coefficient metric whose
curvature dominates a positive base form and whose total-space
multiplier ideal is trivial.  Singular adjoint direct-image positivity
of P\u aun--Takayama~\cite{PT2018} then gives a metric with a global
strict lower bound on $-K_Y+C$.  Its integral boundary has exceptional
pullback, and the distinguished section is locally square integrable.

Keeping the coefficient metric smooth gives a different output.
The coarea formula and a primitive-representative curvature calculation
produce an integral correction $C$ with $\pi_*g^*C=0$ and
\[
 |s_C|^2e^{-\phi}\in L^{1+\eta}_{\mathrm{loc}}
 \quad\text{for some }\eta>0.
\]
The metric is smooth and strictly positive on one ordinary open set.
The higher integrability exponent permits a mixture with a globally
strict metric while keeping $s_C$ integrable.  Thus the global bound
needed for Theorem~\ref{thm:controlled-boundary} follows from an
explicit additional step.

The companion \emph{Metric descent and rank-preserving contractions}
supplies the toroidal-volume construction used here
\cite[Theorem~2.1]{CompanionG}.  For a normal
projective model $p:W\to X$ and a surjective equidimensional toroidal
map $f:W\to Y$, its theorem assumes adjoint rank one at degree zero
on a smooth resolution and a patch, away from the Jacobian divisor,
where the pulled-back curvature dominates a positive base form.
It produces an effective rational correction $A_0$ with
\[
 p_*f^*A_0=0
\]
and a semipositive metric on $-K_Y+A_0$ with unweighted locally
integrable squared frame norms, smooth and strictly positive on a
nonempty ordinary open set.  That companion proves the toroidal
Jacobian calculation, the all-valuation integrability and the
boundary extension.  We use its metric after checking the stated
model, rank and curvature premises in each application.
The rational correction is used directly in
Corollary~\ref{conv:unweighted-effectivity}; no rounding or replacement
of its unweighted metric by an integral-boundary hypothesis is needed.

For the invariant problem, M\"uller \cite[Theorem~C]{Muller2025}
obtains naturally invariant sections for projective sub-log-canonical
pairs with semiample anti-log-canonical divisor and a commutative linear
algebraic group action.  Theorem~\ref{inv:nonvanishing} here assumes
$X$ smooth connected projective, $-K_X$ smoothly semipositive and
$\chi(X,\mathcal O_X)\ne0$.  For any algebraic torus acting on $X$,
it gives
\[
 H^0(X,-mK_X)^T\ne0\quad\text{for some }m>0
\]
with the natural tangent-determinant linearization.
The companion invariant-index theorem, including its actual value at
zero, supplies invariant cohomology
\cite[Theorem~1.1]{CompanionI}; compact-averaged hard Lefschetz
then gives forms with the coefficient shift $m+1$.
The local determinant and character arguments supply the sections
needed here.  The finite-cover structure and norm used in the
isometric application are those of the companion headline paper,
\cite[Theorem~2.1 and Lemma~2.2]{CompanionH}, which develops the geometric work of Demailly--Peternell--Schneider
and Campana--Demailly--Peternell~\cite{DPS1996,CDP2015}.
Our algebraic boundary criterion, smooth coarea proof and correction
ring are independent of that application.

Finally, a suitable auxiliary base has an effective big klt boundary
$\Delta$ with $K_Y+\Delta\sim_{\mathbb Q}C$ for an effective rational
$C$ supported on divisors disappearing over $X$.  When $\kappa(Y,C)=0$, a log
terminal model of the base is of Fano type: it admits an effective
rational boundary $\Theta$ with klt singularities such that
$-(K+\Theta)$ is ample.  This is a statement about
the auxiliary model.  Its proof retains the precise contracted-prime
support needed to return the resulting divisors to $X$.

The main route is the boundary and ring results followed by the
single-Iitaka conversion in Section~\ref{conv:section}.
Sections~\ref{coa:analytic-tools} and~\ref{jet:iitaka-section}
retain the smooth coarea estimate and direct jet-separation proof
of bigness; their applications reuse the common return of sections.
Section~\ref{fano:section} then develops the supported Fano-type
model.  The invariant base constructions in Sections~\ref{inv:section}
and~\ref{inv:dominated-alternative} lead to compact-monodromy descent
in Section~\ref{inv:isometric-descent}.
Finally, Section~\ref{ext:section} constructs invariant tensors from
extremal characters of the full equivariant Euler characteristic.
That construction is independent of the weight-zero index input and
descends tensors before applying ordinary conversion.

\section{From an integrable section to a rational boundary}
\label{sec:approximation}

The analytic information in Theorem~\ref{thm:controlled-boundary}
concerns one section, rather than a nonvanishing local frame.  The
distinction already appears on a disk.  For a weight
$\phi=a\log|z|^2$ with $a\ge0$ and a section $s_J=z^j$ with
$j\in\mathbb Z_{\ge0}$, the two integrability thresholds are
\[
 |s_J|^2e^{-\phi}\in L^1_{\mathrm{loc}}
       \quad\Longleftrightarrow\quad a<j+1,
 \qquad
 e^{-\phi}\in L^1_{\mathrm{loc}}
       \quad\Longleftrightarrow\quad a<1.
\]
On a resolution, the Jacobian contributes an additional vanishing
order.  After dividing the fixed divisor by the multiple used for
approximation, subtracting this Jacobian order and that of $J$ will
leave coefficients strictly below one.  The resulting signed divisor
has a positive part that can enter a klt boundary, while its negative
part records the
support on which a correction is permitted.

We use additive notation for rational line bundles.  If a metric has
local squared frame norm $e^{-\phi}$, its curvature is positive when
$\phi$ is plurisubharmonic.  A global strict lower bound means
$\sqrt{-1}\partial\bar\partial\phi\ge c\omega$ for one K\"ahler form
$\omega$ and one $c>0$.  Smooth changes of weight do not affect local
integrability.  The multiplier ideal $\mathcal J(\phi)$ consists of
holomorphic germs $u$ for which $|u|^2e^{-\phi}$ is locally integrable.
Section~\ref{sec:correction-ring} uses the unweighted case $J=0$;
it first obtains a global strict bound while preserving integrability
of squared frame norms.

\begin{lemma}[Approximation retaining a distinguished section]
\label{lem:weighted-resolution}
Let $Y$ be smooth projective, let $J\ge0$ be an integral divisor, and
let $B$ be a rational line bundle with a singular metric $h$ having a
global strict curvature lower bound.  Suppose
$\mathcal J(h)\supset\mathcal O_Y(-J)$.
For every sufficiently large divisible integer $r$, the complete system
$|rB|$ is nonempty.  Let $\sigma:Y'\to Y$ be a log resolution of its
base ideal, $J$, and the exceptional locus, and write
\[
 \sigma^*|rB|=G+|\mathrm{Mov}|,
\]
where $G$ is the fixed divisor and $|\mathrm{Mov}|$ is base-point free.
Then the simple-normal-crossing divisor
\begin{equation}\label{eq:weighted-signed-divisor}
 \Psi=G/r-\sigma^*J-K_{Y'/Y}
\end{equation}
has every coefficient strictly less than one.  Its negative part is
supported on $\sigma^*J$ and the exceptional primes.
\end{lemma}

\begin{proof}
Fix a very ample integral divisor $H$ and put $d=\dim Y$.  For $r$
sufficiently large and divisible, the curvature of $rh$ dominates the
fixed smooth curvatures of $K_Y+iH$, simultaneously for $1\le i\le d$.
Nadel vanishing~\cite{Nadel1990,Demailly2015} gives
\[
 H^i\bigl(Y,\mathcal O_Y(rB-iH)\otimes\mathcal J(r\phi)\bigr)=0
 \quad (1\le i\le d).
\]
Indeed, in adjoint form the coefficient line is $rB-K_Y-iH$, whose
metric has a strict positive lower bound.  Castelnuovo--Mumford
regularity makes
$\mathcal O_Y(rB)\otimes\mathcal J(r\phi)$ globally generated.
This is a nonzero coherent sheaf of rank one.  Consequently $|rB|$ is
nonempty and its base ideal $\mathfrak b_r$ satisfies
\begin{equation}\label{eq:base-ideal-multiplier}
 \mathfrak b_r\supset\mathcal J(r\phi).
\end{equation}

We next compare the fixed divisor with the singularity of $h$.  Fix $r$
and a prime $S$ on $Y'$.  At its general point choose a transverse
coordinate $t$ and write
\[
 a=\ord_SG,\qquad j=\ord_S\sigma^*J,
 \qquad k=\ord_SK_{Y'/Y}.
\]
On nested coordinate balls downstairs, the point-extension estimate
of Ohsawa--Takegoshi, in its weighted Bergman formulation
\cite{OhsawaTakegoshi1987,Demailly2015}, gives holomorphic functions
$u_z$ with
\[
 |u_z(z)|^2=e^{r\phi(z)},\qquad
 \int |u_z|^2e^{-r\phi}\le C_r
\]
for every center $z$ in the smaller ball where $\phi(z)>-\infty$.
The constant is independent of $z$; it need not be independent of $r$.
Since $\phi$ is bounded above on a slightly smaller ball, holomorphic
mean-value estimates bound $u_z$ uniformly there.  Each germ of $u_z$
belongs to $\mathcal J(r\phi)$, hence to $\mathfrak b_r$ by
\eqref{eq:base-ideal-multiplier}.  Its pullback is therefore divisible
by $t^a$ near a general point of $S$.  Cauchy estimates on nested
polydisks bound $(u_z\circ\sigma)/t^a$ uniformly as well.  Substituting
$z=\sigma(w)$ proves
\begin{equation}\label{eq:weighted-local-order}
 \phi\circ\sigma\le \frac ar\log|t|^2+O(1).
\end{equation}
The assertion is automatic on the polar set of $\phi$.

The density $|s_J|^2e^{-\phi}$ is locally integrable downstairs.
Change of variables contributes a factor of order $2k$ along $S$.
By \eqref{eq:weighted-local-order}, its pullback is bounded below,
up to a positive bounded factor, by
$|t|^{2(j+k-a/r)}$ at a general point of $S$.  Integrability in the
transverse complex variable forces
\[
 \frac ar-j-k<1.
\]
These are exactly the coefficients in
\eqref{eq:weighted-signed-divisor}.  Resolving the three supports makes
that divisor simple normal crossings.  Away from $\sigma^*J$ and the
exceptional locus, its only possible contribution is $G/r\ge0$,
which proves the assertion about its negative part.
\end{proof}

The strict inequalities are the essential output.  A non-strict
inequality would not produce a klt boundary.  The local extension
argument explains why integrability of the distinguished section is
precisely the hypothesis that survives after subtracting $\sigma^*J$.

\begin{theorem}[Weighted anticanonical approximation]
\label{thm:weighted-approximation}
Let $Y$ be smooth projective and let $J\ge0$ be integral.  Suppose
$-K_Y+J$ has a singular metric with a global strict curvature lower
bound and multiplier ideal containing $\mathcal O_Y(-J)$.
Choose a sufficiently small ample rational divisor $A$.
For a sufficiently large divisible integer $r$, resolve the complete
system $|r(-K_Y+J-A)|$ and $J$ by $\sigma:Y'\to Y$, and write
\[
 \sigma^*|r(-K_Y+J-A)|=G+|\mathrm{Mov}|.
\]
Then $\Psi=G/r-\sigma^*J-K_{Y'/Y}$ has coefficients less than one,
$\Psi^+$ is an effective klt boundary, and
$D_0=\Psi^-$ is effective and supported on $\sigma^*J$ and exceptional
primes.  Moreover
\begin{equation}\label{eq:weighted-adjoint-identity}
 D_0\sim_{\mathbb Q}
 K_{Y'}+\Psi^++\mathrm{Mov}/r+\sigma^*A.
\end{equation}
The rational class $\mathrm{Mov}/r+\sigma^*A$ is nef and big.
If an algebraic torus acts on $Y$, $J$ is invariant, and the indicated
full systems are torus-stable, the resolution may be chosen equivariantly.  Then $G$,
$K_{Y'/Y}$, $\Psi^\pm$ and $D_0$ are invariant, and the free
system $|\mathrm{Mov}|$ is torus-stable.  Identity
\eqref{eq:weighted-adjoint-identity} is an identity of ordinary rational
line bundles; it does not assert that a chosen member of either the
free system or the ample class is invariant.
\end{theorem}

\begin{proof}
Subtracting a sufficiently small smooth ample metric preserves the
strict lower bound and does not change the multiplier ideal.
Apply Lemma~\ref{lem:weighted-resolution} to $B=-K_Y+J-A$.
Its identity $G+\mathrm{Mov}\sim_{\mathbb Q}r\sigma^*B$ gives
\eqref{eq:weighted-adjoint-identity}.  The positive part of a simple
normal crossing divisor with coefficients less than one is klt.
The moving class is nef, and the pullback of the ample class $A$ is
nef and big, proving the last positivity assertion.
In the equivariant setting the complete-system base ideal is invariant;
equivariant resolution in characteristic zero, applied also to $J$,
therefore supplies invariant fixed, relative canonical and boundary
divisors.  This does not require the original metric to be invariant.
\end{proof}

The signed identity supplies more information than an effectivity
conclusion.  In particular, its negative part has an explicit support.
It does not assert integrability of the norm of a nonvanishing frame;
that stronger hypothesis enters Section~\ref{sec:correction-ring}.

\section{Effectivity with a prescribed correction support}
\label{sec:effectivity}

The signed approximation separates two kinds of data: the negative
part is confined to the permitted correction support, while the ample
part can absorb a small multiple of any prescribed rational divisor.
We use that ample part to insert the pseudoeffective divisor of
Theorem~\ref{thm:controlled-boundary}.  The resulting boundary is
effective, big and klt, so big-boundary nonvanishing applies to its
adjoint.  We first record the passage from its real-linear conclusion
to the rational linear equivalence required here.

\begin{lemma}[Rational recovery]\label{lem:rational-recovery}
Let $D$ be a rational divisor on a normal projective variety.
If $D$ is real-linearly equivalent to an effective real divisor, then
it is rationally linearly equivalent to an effective rational divisor.
\end{lemma}
\begin{proof}
Write an effective real representative as
$D+\sum_{j=1}^s a_j\operatorname{div}(f_j)$, with $a_j\in\mathbb R$
and rational functions $f_j$.  The supports of $D$ and of these finitely
many principal divisors give a finite system of rational linear
inequalities asserting nonnegativity of every coefficient.
The known real solution makes the resulting rational polyhedron
nonempty.  Such a polyhedron has a rational point: take a point in the
relative interior of its smallest supporting face, solve that face's
rational linear equations over $\mathbb Q$, and approximate within
its rational affine space while preserving the remaining strict
inequalities.  Substituting that rational point gives the assertion.
\end{proof}

\begin{proof}[Proof of Theorem~\ref{thm:controlled-boundary}]
Fix a small ample rational divisor $A$ whose smooth metric can be
subtracted from the given metric on $-K_Y+C$ while retaining its global
strict curvature lower bound.  Apply
Theorem~\ref{thm:weighted-approximation} with $J=C$ and choose $r>1$.
On the resulting resolution $\sigma:Y'\to Y$, write
$\Psi=\Psi^+-\Psi^-$ and retain the free system $|\mathrm{Mov}|$.
The signed adjoint identity is
\[
 K_{Y'}+\Psi^++\mathrm{Mov}/r+\sigma^*A
       \sim_{\mathbb Q}\Psi^-.
\]
Choose a general $D_{\mathrm{mov}}\in|\mathrm{Mov}|$ meeting the
resolved support transversely and without common components.
All these data are fixed before choosing the pseudoeffective divisor $M$.

The ample cone is open, so $A+tM$ is ample for a sufficiently small
positive rational $t$.  Choose an effective rational representative
$U_t\sim_{\mathbb Q}A+tM$ by dividing a general member of a sufficiently
high very ample multiple.  Bertini permits $\sigma^*U_t$ and $D_{\mathrm{mov}}$ to meet
the resolved support with simple normal crossings, with no common
components; the new coefficients are less than one.  Thus
\begin{equation}\label{eq:effective-big-boundary}
 \Delta_t=\Psi^++D_{\mathrm{mov}}/r+\sigma^*U_t
\end{equation}
is effective and klt.  Its class is big because it contains the
effective summand $\sigma^*U_t$, the pullback of an ample rational class.
The adjoint identity becomes
\begin{equation}\label{eq:effectivity-adjoint}
 K_{Y'}+\Delta_t\sim_{\mathbb Q}\Psi^-+t\sigma^*M.
\end{equation}
Its right-hand side is pseudoeffective.

Big-boundary nonvanishing of Birkar--Cascini--Hacon--McKernan
\cite[Theorem D]{BCHM2010} therefore gives an effective real divisor
real-linearly equivalent to $K_{Y'}+\Delta_t$.
Lemma~\ref{lem:rational-recovery} makes this a rational linear
equivalence with an effective rational divisor.  Finally,
$F=\sigma_*\Psi^-$ is effective and supported on $C$ by
Lemma~\ref{lem:weighted-resolution}.  Push
\eqref{eq:effectivity-adjoint} down to $Y$ and divide by $t$.
The theorem follows with $C_0=F/t$.
\end{proof}

\begin{proposition}[A fixed support correction before scaling]
\label{prop:ample-perturbation}
Under the hypotheses of Theorem~\ref{thm:controlled-boundary}, there
is an effective rational divisor $F$ supported on $C$ such that for
every pseudoeffective rational divisor $M$, some positive rational
$t$ makes $F+tM$ rationally linearly equivalent to an effective
rational divisor.
\end{proposition}
\begin{proof}
In the preceding proof, the approximation and hence
$F=\sigma_*\Psi^-$ were chosen independently of $M$.  Only $t$ and the
general ample representative depended on $M$.
\end{proof}

When $C=0$, the correction vanishes.  For nonzero $C$, it is the
integrability of $s_C$ that controls the support of $\Psi^-$; no
unweighted frame-integrability hypothesis has been imposed.
The divisor to which nonvanishing is applied is the effective boundary
\eqref{eq:effective-big-boundary}.

\subsection{Preserving integrability under a pseudoeffective perturbation}

One can also insert $M$ at the level of metrics, before approximation.
This gives analytic information in addition to the effectivity proved
above: a small pseudoeffective perturbation preserves the integrability
of the distinguished section and the global strict curvature bound.

\begin{lemma}\label{lem:weighted-pseudoeffective-perturbation}
Under the hypotheses of Theorem~\ref{thm:controlled-boundary}, for
every pseudoeffective rational divisor $M$ there is a positive rational
$\epsilon$ such that $-K_Y+C+\epsilon M$ has a metric with a global
strict curvature lower bound and multiplier ideal containing
$\mathcal O_Y(-C)$.
\end{lemma}
\begin{proof}
Write $\phi$ for the given metric and choose a semipositive singular
metric with weights $\theta$ on $M$.  Strong openness
\cite{GuanZhou2015} gives a common $\tau>0$ on a finite relatively
compact coordinate covering such that
$|s_C|^2e^{-(1+\tau)\phi}$ is locally integrable.  Put
$q=(1+\tau)/\tau$.  H\"older's inequality with respect to the measure
$|s_C|^2d\lambda$ gives
\[
 \int |s_C|^2e^{-\phi-\epsilon\theta}
 \le
 \left(\int |s_C|^2e^{-(1+\tau)\phi}\right)^{1/(1+\tau)}
 \left(\int |s_C|^2e^{-q\epsilon\theta}\right)^{1/q}.
\]
Small-exponent integrability of plurisubharmonic functions
\cite{Skoda1972} makes the second factor finite for sufficiently small
positive rational $\epsilon$.  The finite covering permits one choice
of $\epsilon$.  The weight $\phi+\epsilon\theta$ retains the global
strict curvature lower bound because the metric on $M$ is semipositive.
\end{proof}

Subtract a sufficiently small ample rational $A$ from the perturbed
metric, retaining its strict lower bound.  Applying
Lemma~\ref{lem:weighted-resolution} to $-K_Y+C+\epsilon M-A$ with
$r>1$ then gives a second boundary construction.
For its signed divisor $\Psi$ and a general free member
$D_{\mathrm{mov}}$, add a general small-coefficient representative
$U\sim_{\mathbb Q}A$.  Then
$\Delta=\Psi^++D_{\mathrm{mov}}/r+\sigma^*U$ is effective, big and klt,
and $K_{Y'}+\Delta\sim_{\mathbb Q}\epsilon\sigma^*M+\Psi^-$.
Nonvanishing, rational recovery and pushforward apply as before.
Here the approximation itself depends on the metric on $M$;
the support conclusion is unchanged.

\section{Unweighted integrability and nearby adjoint rings}
\label{sec:correction-ring}

The criterion just proved accepts a weighted integral and a global
strict curvature bound.  For finite generation we instead begin with
integrable squared \emph{frame} norms and strict positivity on one
ordinary open set.  This permits perturbation in arbitrary finitely
many rational directions.  We first make the passage from local
strictness to a global lower bound while preserving integrability.

\begin{lemma}[An integrable neighborhood]\label{lem:integrable-neighborhood}
Let $Y$ be smooth projective of positive dimension and let $B$ be a
rational line bundle.  Suppose $B$ has a semipositive singular metric
with locally integrable squared frame norms, smooth and strictly
positive on a nonempty ordinary open set.
For any finite list of rational line bundles $P_1,\ldots,P_s$, there
are $\epsilon\in\mathbb Q_{>0}$ and a common ample rational line
bundle $H'$ such that each $B+\epsilon P_i-H'$ is rationally linearly
equivalent to an effective divisor $\Delta_i$ with $(Y,\Delta_i)$ klt.
\end{lemma}

\begin{proof}
The absolutely continuous part of the curvature current has positive
top-degree mass on the given open set.  Boucksom's volume criterion
\cite[Theorem 1.2(i)]{Boucksom2002} implies that $B$ is big.
A rational big decomposition therefore gives a metric with algebraic
singularities and weights $\psi$ whose curvature dominates $c\omega$
for some $c>0$.  Write $\phi$ for the original weights.
Strong openness and compactness give $\eta>0$ with
$e^{-(1+\eta)\phi}$ locally integrable on one finite covering.
Small-exponent integrability supplies $a>0$ with $e^{-a\psi}$
locally integrable on that covering as well.

Choose $p>1$ with $p<1+\eta$ and let $q=p/(p-1)$.
For small positive rational $u$ we have $p(1-u)<1+\eta$ and $qu<a$.
H\"older's inequality then makes
\[
 e^{-((1-u)\phi+u\psi)}
\]
locally integrable.  The mixed metric has curvature at least $uc\omega$.
This is the required global strict lower bound; positivity on the
original open set was used only to establish bigness.

Choose smooth metrics on all $P_i$ and on one ample integral divisor
$H$.  Their curvatures have uniformly bounded norm relative to $\omega$.
Taking rational $\epsilon>0$ and $h>0$ sufficiently small, subtracting
$H'=hH$ and adding $\epsilon P_i$ leaves every mixed metric with a
common strict lower bound.  Smooth changes of weights preserve
integrability.

It remains to turn a rational line bundle $D$ with such a strict
integrable metric into an effective klt divisor.  Apply
Lemma~\ref{lem:weighted-resolution} with $J=0$ and $r>1$ to $D$.
Take a general member of $|rD|$ and denote its divisor divided by $r$
by $\Delta$.  On the resolution its pullback is
$G/r+D_{\mathrm{mov}}/r$.  The lemma says that
$G/r-K_{Y'/Y}$ has coefficients less than one.
A general free member meets the fixed support transversely and has
coefficient $1/r<1$.  Thus all coefficients of
$\sigma^*\Delta-K_{Y'/Y}$ are less than one with simple normal
crossing support.  This is exactly the klt condition for the effective
pair $(Y,\Delta)$.  Since $\Delta\sim_{\mathbb Q}D$, apply this
argument to $D=B+\epsilon P_i-H'$ for each $i$.
\end{proof}

\begin{corollary}\label{cor:big-klt-representative}
Under the hypotheses of Lemma~\ref{lem:integrable-neighborhood}, $B$
has an effective big klt rational representative.  The same holds for
$B+\epsilon P_i$ for the sufficiently small simultaneous perturbations
in that lemma.
\end{corollary}

\begin{proof}
Add to $\Delta_i$ a general effective rational representative of $H'$
with sufficiently small coefficients.  On a common log resolution,
Bertini preserves the klt inequalities.  The sum is big because it
contains an effective representative of the ample class $H'$.
For $B$ itself, take the sole direction $P_1=0$.
\end{proof}

\begin{theorem}[Finite generation near a correction divisor]
\label{thm:correction-ring}
Let $Y$ be smooth projective of positive dimension, and let $A_0$ be
a rational divisor.  Suppose $-K_Y+A_0$ has a semipositive singular
metric with locally integrable squared frame norms, smooth and
strictly positive on a nonempty ordinary open set.
Given rational divisors $P_1,\ldots,P_s$, there are
$\epsilon\in\mathbb Q_{>0}$ and an integer $k_0>0$ such that the
actual divisors $B_i=k_0(A_0+\epsilon P_i)$ are Cartier and
\[
 R(Y;B_1,\ldots,B_s)=
 \bigoplus_{(a_1,\ldots,a_s)\in\mathbb N^s}
 H^0\!\left(Y,\mathcal O_Y\Bigl(\sum_i a_iB_i\Bigr)\right)
\]
is a finitely generated $\mathbb C$-algebra.
\end{theorem}

\begin{proof}
Lemma~\ref{lem:integrable-neighborhood} gives effective klt rational
boundaries $\Delta_i$ and one ample rational $H'$ such that
\[
 \Delta_i\sim_{\mathbb Q}-K_Y+A_0+\epsilon P_i-H'.
\]
The multigraded adjoint finite-generation theorem of
Corti--Lazi\'c~\cite[Theorem 3.2(1)]{CortiLazic2013}, based on
Cascini--Lazi\'c~\cite{CasciniLazic2012}, applies to the divisors
$K_Y+\Delta_i+H'$.  Clear all denominators and all rational linear
equivalences with a common integer $k_0$.  Multiplication by the
corresponding rational functions, compatibly in each multidegree,
identifies a multigraded Veronese of that adjoint ring with the displayed
ring.  Finite-index multigraded Veronese subrings are finitely generated
\cite[Lemma 3.1]{CortiLazic2013}.
The construction uses rational linear equivalence and actual Cartier
divisors, not only their numerical classes.
\end{proof}

\subsection{Direct selection of an integrable divisor}
\label{subsec:bergman-selection}

There is a second way to perform the last step of
Lemma~\ref{lem:integrable-neighborhood}.  It chooses a section by
averaging negative powers of its norm, instead of using a general
member on a log resolution.  This gives a direct analytic description
of the effective klt representative.

Let $D$ be a rational line bundle with a globally strictly positive
metric of weights $\phi_D$, and suppose $e^{-\phi_D}$ is locally
integrable.  For a large divisible integer $m>1$, take an orthonormal
basis $s_1,\ldots,s_N$ of the global sections of $mD$ that are square
integrable with weight $m\phi_D$ and a fixed smooth volume form.
The global Bergman estimate is
\begin{equation}\label{eq:global-bergman-selection}
 \frac1m\log\sum_{j=1}^N|s_j|^2\ge\phi_D-C_m
\end{equation}
on each relatively compact coordinate patch.  We recall the
construction because this version uses global sections of $mD$
without a fixed auxiliary twist.

Fix a point $y$ where $\phi_D(y)$ is finite.  Point extension on a
coordinate ball gives a local holomorphic section taking squared
value $e^{m\phi_D(y)}$ at $y$, with uniformly bounded weighted norm.
Multiply it by a cutoff equal to one on a smaller ball and solve the
resulting $\bar\partial$ equation.  In top holomorphic degree the
coefficient line is $mD-K_Y$.  Add a cutoff logarithmic pole
$d\log|z-y|^2$ to the estimating weight, where $d=\dim Y$.
For sufficiently large $m$, the strict curvature of $mD$ absorbs both
the fixed canonical curvature and the negative curvature of the cutoff
pole, uniformly for $y$ in a smaller ball.  The correction is
holomorphic near $y$; its integrability against $|z-y|^{-2d}$ forces
its value there to vanish.  The corrected global section retains the
specified value and has uniformly bounded $m\phi_D$-weighted norm.

For the singular-weight solution, first choose rational trivializing
sections of an integral multiple of $D$ and of $K_Y$.  Enlarge the
support of their zeros and poles by an effective divisor so that the
resulting divisor is ample.  Its complement is affine and trivializes
the required integral line bundles.  Use Stein exhaustions of that
complement, approximate the psh weights on relatively compact
exhaustion sets and pass to weak limits.  In bidegree $(d,1)$ the
$L^2$ estimate uses contraction with the inverse positive
line-curvature matrix; its bound is independent of an auxiliary
complete K\"ahler metric chosen to dominate a fixed reference metric.
The estimates therefore persist through the exhaustion.  Since the
original psh weights are locally bounded above, weighted square
integrability implies ordinary square integrability near the removed
divisor.  The resulting holomorphic section extends across it.
The local weighted estimate is Demailly's Bergman approximation
\cite[Section~1, Theorem~1.2]{Demailly2015}, with the point-extension
input of Ohsawa--Takegoshi~\cite{OhsawaTakegoshi1987}.  The strict
global curvature is what permits the untwisted global construction
just given.
Expanding the constructed section in the orthonormal basis proves
\eqref{eq:global-bergman-selection}; the estimate is automatic on the
polar set of the weight.

It follows that $(\sum_j|s_j|^2)^{-1/m}$ is locally integrable.
Average over the unit sphere in $\mathbb C^N$.  Unitary invariance gives
\[
 \int_{\|a\|=1}\left|\sum_{j=1}^Na_js_j(y)\right|^{-2/m}\,da
 =c_{m,N}\left(\sum_{j=1}^N|s_j(y)|^2\right)^{-1/m}.
\]
For $N\ge2$, the possible singularity is an inverse power of one
complex coordinate and $c_{m,N}<\infty$ because $m>1$.
For $N=1$ the identity is immediate.
Fubini on a finite covering gives a nonzero section $s$ of $mD$
with $|s|^{-2/m}$ locally integrable.  On a log resolution, change of
variables identifies this condition with all discrepancy coefficients
of $\operatorname{div}(s)/m$ being less than one.  Thus
$\operatorname{div}(s)/m$ is the required effective klt rational
representative of $D$.

\section{Removing a fixed pseudoeffective error}
\label{conv:section}

The controlled-boundary theorem becomes useful for anticanonical
nonvanishing when the boundary on an auxiliary base disappears after
pullback and birational pushforward.  We now construct precisely such a
base.  The error in the following theorem is an arbitrary fixed
pseudoeffective divisor; it need not be effective.

\begin{theorem}\label{conv:fixed-error}
Let $X$ be a smooth connected projective complex variety and let
$L=-K_X$ have a smooth semipositive Hermitian metric.  Let $P$ be a
pseudoeffective Cartier divisor.  If there are effective integral
$N_j\sim m_jL-P$ for unbounded positive integers $m_j$, then
$H^0(X,mL)\ne0$ for some positive integer $m$.
\end{theorem}

We use the Iitaka fibration of just one $N_j$.  Ratios of its sections
force all subsequent divisors to vary affinely on the generic fiber.
The vertical slopes define a pseudoeffective divisor on the base.  A
rank-one adjoint direct image supplies the metric required by
Theorem~\ref{thm:controlled-boundary}, and its boundary is exceptional
over $X$.  Thus the correction allowed by that theorem causes no poles
in the section returned to $X$.

\subsection{Vertical minima and return of sections}

We record the return step once, including the exceptional divisors that
can appear when a smooth model ceases to be equidimensional.

\begin{lemma}\label{conv:return}
Let $\pi:Z\to X$ be a projective birational morphism of smooth
projective varieties and let $f:Z\to Y$ be a surjective morphism to a
smooth projective variety with connected general fibers, and let $L$ be
a rational line bundle on $X$.  Assume every
prime of $Z$ mapping into codimension at least two in $Y$ is
$\pi$-exceptional.  Let $R\sim_{\mathbb Q}a\pi^*L$, $a>0$, be a
rational divisor with nonnegative horizontal part.  For each base prime
$Q$ put
\[
 \mu_Q(R)=\min_{D\mapsto Q}
 \frac{\operatorname{ord}_D R}{\operatorname{ord}_D f^*Q},
 \qquad M=\sum_Q\mu_Q(R)Q.
\]
Suppose $M$ is pseudoeffective.  If there is an effective rational
$C_0$, all components of whose pullback are $\pi$-exceptional, such that
$M+C_0\sim_{\mathbb Q}G\ge0$, then $L$ is rationally linearly
equivalent to an effective rational divisor.  The same conclusion holds
if the premise is $\epsilon M+C_0\sim_{\mathbb Q}G\ge0$ for some
rational $\epsilon>0$.

For the invariant conclusion, give $L$ a specified linearization and
assume
\[
 R=c\,\operatorname{div}(\pi^*t),\qquad cr=a,
\]
where $c>0$ is rational and $t$ is a nonzero invariant rational section
of the specified linearized $L^r$, with $r>0$ an integer for which
$L^r$ is a line bundle.  Assume also that every function pulled back
from $Y$ is invariant in $\mathbb C(X)$.  Then the resulting section
of a positive multiple of $L$ is invariant for that linearization;
no regular group action on $Z$ is required.
\end{lemma}
\begin{proof}
Only finitely many coefficients of $M$ are nonzero.  Replace $C_0,G$ by
$C_0/\epsilon,G/\epsilon$ in the last variant.  The divisor
\[
 R+f^*(G-M-C_0)\sim_{\mathbb Q}a\pi^*L
\]
has nonnegative horizontal coefficients.  At a prime dominating $Q$
outside the support of $C_0$, this follows from the minimum defining
$M$ and from $G\ge0$.  All other possibly negative primes are
$\pi$-exceptional, either by the support premise on $C_0$ or by the
assumption on primes mapping into codimension two.  Its pushforward is
therefore effective and rationally linearly equivalent to $aL$.
Clearing denominators proves the assertion.  In the invariant case
choose the clearing integer $n$ also so that $nc$ is integral.  The
operation multiplies $t^{nc}$, a section of $L^{ncr}=L^{na}$, by a
rational function pulled from $Y$.  Both factors are invariant, and
extension across codimension two on the smooth variety $X$ preserves
invariance.
\end{proof}

The required model condition follows from flattening, rather than from
an assertion that a final smooth resolution is flat.  Starting with a
resolved rational map and its Stein factorization, flatten over a
modification of the base \cite[Theorem~5.2.2]{RaynaudGruson1971}.
Normalize the main component, resolve the base, flatten again if
necessary, and then resolve the source.  On the intervening
flat model a prime cannot map into base codimension two: its inverse
image has dimension at most $\dim X-2$.  Any prime on the final
resolution with that property is consequently exceptional over the
intervening model, and hence over $X$.

\subsection{The Iitaka base and its strict metric}
\label{conv:iitaka}

\begin{proof}[Proof of Theorem~\ref{conv:fixed-error}]
Pass to a subsequence with $m_0<m_1<m_2<\cdots$, and take the Iitaka
fibration of $N_1$.  Resolve a system of maximal image dimension, take
its Stein factorization, and prepare the model by flattening as above:
\[
 X\xleftarrow{\ \pi\ }Z\xrightarrow{\ f\ }Y.
\]
The base field $\mathbb C(Y)$ is relatively algebraically closed in
$\mathbb C(X)$ and contains every ratio of sections of every positive
multiple of $N_1$.  Indeed, a ratio transcendental over the original
image field would increase the image dimension after multiplication of
sections into a common system.  All other ratios lie in its relative
algebraic closure, which is the Stein field.  If $\dim Y>0$, we may
choose an ample $H$ on $Y$ and an integer $b>0$ with
\begin{equation}\label{conv:iitaka-polarization}
 b\pi^*N_1\sim f^*H+D_1,\qquad D_1\ge0.
\end{equation}
The polarization of the original image pulls back to a big divisor on
$Y$; a sufficiently large multiple dominates $H$, which proves this
assertion on the modified base.

Put $R=(N_1-N_0)/(m_1-m_0)\sim_{\mathbb Q}L$.  For $j>1$ the two
effective divisors
\[
 (m_j-m_1)N_0+(m_1-m_0)N_j,
 \qquad (m_j-m_0)N_1
\]
are linearly equivalent.  Their ratio is a base function.  Consequently
there are rational principal divisors $Q_j\sim_{\mathbb Q}0$ on $Y$
with the exact identities
\begin{equation}\label{conv:affine-divisors}
 \pi^*N_j=\pi^*N_0+(m_j-m_0)(\pi^*R+f^*Q_j).
\end{equation}
Effectivity as $m_j\to\infty$ makes the horizontal part of $\pi^*R$
nonnegative.  If $Y$ is a point, this already proves the theorem.  If
$\dim Y=\dim X$, then $N_1$ is big, so $m_1L\sim N_1+P$ is big and
again has a section in a positive multiple.  We henceforth assume
$0<\dim Y<\dim X$.

Define $M=\sum_Q\mu_Q(\pi^*R)Q$ as in
Lemma~\ref{conv:return}.  There is a fixed effective rational $J$ on
$Y$ such that
\[
 M+Q_j+\frac{J}{m_j-m_0}\ge0.
\]
For example its coefficient at $Q$ can be the maximum of
$\operatorname{ord}_D\pi^*N_0/\operatorname{ord}_D f^*Q$ over
$D\mapsto Q$.  These maxima have finite support.  The inequality
follows by taking each coefficient in \eqref{conv:affine-divisors}.
Since $Q_j$ is rational principal, the numerical classes of these
effective divisors converge to $[M]$.  Thus $M$ is pseudoeffective.

The remaining problem is to make $M$ effective with a correction
exceptional over $X$.  Write $E=K_{Z/X}\ge0$.  We first show
\begin{equation}\label{conv:adjoint-rank}
 \operatorname{rank}f_*\mathcal O_Z(E)=1.
\end{equation}
The Jacobian section supplies rank at least one.  If the rank were
larger, twisting by $kH$ for large $k$ would supply two sections of
$E+kf^*H$ independent on the generic fiber.  Multiply by the section of
$bk\pi^*N_1-kf^*H$ furnished by
\eqref{conv:iitaka-polarization}.  The resulting ratio belongs to
$|E+bk\pi^*N_1|$.  Since
$\pi_*\mathcal O_Z(E)=\mathcal O_X$, it is a ratio from a multiple of
$N_1$, and hence a base function, a contradiction.  The equality of
pushforwards used here follows from normality of $X$: a rational
function with possible poles only on exceptional divisors has no pole
at any prime of $X$ and extends across codimension two.

Write
\[
 (f_*\mathcal O_Z(E))^{**}=\mathcal O_Y(C),\qquad C\ge0,
\]
where $C$ is the divisor of the canonical section.  The divisor $C$ is
integral.  The torsion-free rank-one direct image agrees with its hull
at every codimension-one point.  If $Q\subset\operatorname{Supp}C$,
dividing the canonical section by a local equation of $Q$ is therefore
allowed at its generic point.  Hence every prime of $Z$ dominating
$Q$ occurs in $E$.  The model condition handles components over higher
codimension, so every component of $f^*C$ is exceptional over $X$.

We now verify all analytic premises of
Theorem~\ref{thm:controlled-boundary}.  Combining
\eqref{conv:iitaka-polarization} with $N_1\sim m_1L-P$ gives
\[
 bm_1\pi^*L\sim b\pi^*P+f^*H+D_1.
\]
A positive-current metric for $P$, an ample metric for $H$, and the
divisor metric of $D_1$ therefore give a singular metric on $\pi^*L$
whose curvature dominates a positive multiple of $f^*\omega_Y$.
Mix a sufficiently small positive fraction of its weight with the
smooth semipositive pullback metric.  Skoda's small-exponent
integrability theorem, on a finite cover of the compact $Z$, gives a
metric $h$ satisfying
\begin{equation}\label{conv:strict-upstairs}
 \Theta_h\ge\delta f^*\omega_Y,
 \qquad \mathcal J(h)=\mathcal O_Z,
 \qquad \delta>0.
\end{equation}
The metric $h$ need not be smooth.  The triviality of its multiplier
ideal is a statement on the whole total space; a uniform claim on
all fibers is neither needed nor asserted.

Apply singular adjoint direct-image positivity
\cite[Theorem~3.3.5]{PT2018}; see also
\cite[Theorem~21.1 and Corollary~21.2]{HPS2018}.
The natural inclusion of the multiplier-ideal direct image into the
ordinary adjoint direct image is an equality here, because
$\mathcal J(h)=\mathcal O_Z$ in \eqref{conv:strict-upstairs}.
In particular the generic-isomorphism hypothesis in the cited theorem
holds, and it gives a metric on the ordinary sheaf.  Projection formula and
\eqref{conv:adjoint-rank} identify its rank-one hull as
\[
 \bigl(f_*\mathcal O_Z(K_{Z/Y}+\pi^*L)\bigr)^{**}
   =\mathcal O_Y(-K_Y+C).
\]
It acquires a metric $g$ with $\Theta_g\ge\delta\omega_Y$.
For the lower bound, subtract a local potential of
$\delta f^*\omega_Y$ upstairs; the fiber integral subtracts the
corresponding base potential downstairs.  The direct-image theorem
extends the metric across divisors; plurisubharmonic extension handles
the remaining codimension-two set.  Thus this is a global strict
curvature bound on the actual line bundle, including its divisorial
correction.

The distinguished section has finite norm.  On the smooth base-change
locus its norm times coordinate base volume is the pushforward of the
squared norm of the Jacobian section of $K_Z+\pi^*L$.
The latter is locally integrable by $\mathcal J(h)=\mathcal O_Z$.
Properness and Fubini give finite local mass on the base; the omitted
analytic subsets have measure zero.  In the line $-K_Y+C$ the
distinguished local section is $s_C$.  Consequently
\[
 \mathcal J(g)\supset\mathcal O_Y(-C).
\]
We have verified global strictness, the precise multiplier-ideal
inclusion, and exceptional pullback of the integral boundary.
Theorem~\ref{thm:controlled-boundary} gives
$M+C_0\sim_{\mathbb Q}G\ge0$, with $C_0\ge0$ supported on $C$.
Lemma~\ref{conv:return} now gives the required section on $X$.
\end{proof}

\subsection{The smooth toroidal metric for the same base}
\label{conv:toroidal-alternative}

A toroidal embedding is locally modeled on a toric variety with its
dense torus; its boundary corresponds to the toric boundary.  We use
strict toroidal embeddings, whose boundary components are normal.
A toroidal morphism is locally a toric morphism, described on the dense
tori by monomials.  This local description keeps track of the exact
ramification exponents in the metric correction.

There is another way to supply the metric on the Iitaka base.  It keeps
the original smooth anticanonical metric throughout fiber integration
and gives unweighted integrability.  We give the geometric verification
because the two metric constructions have different conclusions.

\begin{lemma}\label{conv:null-curvature}
Let $Z$ be smooth projective of dimension $n$, and let $\theta,\eta$ be
smooth semipositive closed $(1,1)$-forms.  Suppose $a[\theta]-[\eta]$
is pseudoeffective for some $a>0$.  If $\nu$ is the maximum rank of
$\theta$, then $\ker\theta\subset\ker\eta$ on its rank-$\nu$ locus.
If $\eta=f^*\omega_Y$ on an ordinary open meeting the rank-$\nu$
locus, with $f$ submersive and $\omega_Y>0$, there is a smaller
nonempty patch on which
$\theta\ge c f^*\omega_Y$ for some $c>0$.
\end{lemma}
\begin{proof}
There is nothing to prove when $\nu=n$.  Otherwise choose a K\"ahler
form $\omega$.  Every mixed intersection of $\theta$ and $\eta$ of
total degree $\nu+1$, completed by $\omega^{n-\nu-1}$, is
nonnegative.  Replacing one $\eta$ at a time by $a\theta$ can only
increase the intersection, because a pseudoeffective class intersects
nonnegatively with $n-1$ nef classes.  The resulting upper bound is a
multiple of $\int_Z\theta^{\nu+1}\omega^{n-\nu-1}=0$.
Thus $(\theta+\eta)^{\nu+1}=0$ pointwise, by semipositivity.  On the
rank-$\nu$ locus $\theta+\eta$ has the same kernel as $\theta$;
this is the asserted inclusion.  On a relatively compact constant-rank
patch, the nonzero eigenvalues of $\theta$ have a positive lower bound
and $f^*\omega_Y$ is bounded, giving the final inequality.
\end{proof}

We also isolate the algebraic consequence of unweighted integrability.

\begin{corollary}\label{conv:unweighted-effectivity}
Let $Y$ be smooth projective of positive dimension and $A_0\ge0$
rational.  Suppose $B=-K_Y+A_0$ has a semipositive metric with locally
integrable squared frame norms, smooth and strictly positive on a
nonempty ordinary open.  For every pseudoeffective rational $M$ there
is a rational $\epsilon>0$ such that
$A_0+\epsilon M$ is rationally linearly equivalent to an effective
rational divisor.
\end{corollary}
\begin{proof}
Apply Lemma~\ref{lem:integrable-neighborhood} in the direction $M$.
It gives an ample rational $H'$, a sufficiently small rational
$\epsilon>0$, and an effective klt boundary
$\Delta_0\sim_{\mathbb Q}B+\epsilon M-H'$.
Represent $H'$ by a general effective rational divisor with small
coefficients transverse to a log resolution of $\Delta_0$.
Its addition gives an effective big klt boundary
$\Delta\sim_{\mathbb Q}B+\epsilon M$.  The log canonical divisor
$K_Y+\Delta\sim_{\mathbb Q}A_0+\epsilon M$ is pseudoeffective.
Big-boundary nonvanishing \cite[Theorem~D]{BCHM2010} applies.
Lemma~\ref{lem:rational-recovery} gives rational linear effectivity.
\end{proof}

Prepare the rational Iitaka map birationally as an equidimensional
strict toroidal morphism $f:W\to Y$, with $W$ normal, $Y$ smooth,
and $p:W\to X$ birational.  Such a preparation follows from weak
toroidalization \cite[Theorem~1.1]{ADK2013} and equidimensional
toroidal subdivision \cite[Proposition~4.4]{AK2000}.  Mark the
nonisomorphism locus in the source boundary before these operations.
No finite base alteration or reduced-fiber assertion is used here.
Take a smooth resolution $Z\to W$, with maps $\pi:Z\to X$ and
$g:Z\to Y$.  The resolution itself need not be equidimensional.
The generic rank-one argument \eqref{conv:adjoint-rank} is unchanged.
Likewise the horizontal interpolation and the vertical minima may be
computed on $W$, where every vertical prime dominates a base prime.

Equation~\eqref{conv:iitaka-polarization}, pulled to $Z$, says that a
multiple of $[\Theta_{\pi^*L}]$ dominates the pullback of a positive
base class pseudoeffectively.  Lemma~\ref{conv:null-curvature} gives a
horizontal strictness patch above the smooth base-change locus, away
from the Jacobian divisor.  The Jacobian section is nonzero there.
The model, rank-one, and strict-patch hypotheses have now been
verified for the original smooth coefficient metric.  These are the
hypotheses of the toroidal-volume theorem in
\cite[Theorem~2.1]{CompanionG}.
That theorem supplies a rational $A_0\ge0$ with
$p_*f^*A_0=0$ and the unweighted integrable metric of
Corollary~\ref{conv:unweighted-effectivity}.  Its coefficients are
\[
 \operatorname{ord}_Q A_0=
 \max\left\{0,\min_{D\mapsto Q}
       \left(\frac{1+\operatorname{ord}_D K_{W/X}}
                        {\operatorname{ord}_D f^*Q}-1\right)\right\}.
\]
The all-valuation integrability and extension assertions here are part
of that theorem, rather than consequences of generic rank one alone.
Corollary~\ref{conv:unweighted-effectivity} and
Lemma~\ref{conv:return}, on a common resolution, complete the same
conversion.  The output here is unweighted integrability for a rational
correction: a nonvanishing frame has finite norm.  The smooth coarea
construction in Section~\ref{coa:analytic-tools} instead gives an
integral correction and an $L^{1+\eta}$ bound for the distinguished
section.  Lemma~\ref{coa:weighted-effectivity} converts that output
to the same supported effectivity premise in
Lemma~\ref{conv:return}.

\subsection{The field of all pseudoeffectively dominated systems}
\label{conv:all-dominated}

A second base construction uses every effective system dominated by a
multiple of $L$.  It is useful when no distinguished divisor should be
chosen in advance.  We retain the different maximality argument and the
explicit integral correction it produces.

Call a subsystem of an integral divisor $D$ admissible if it has a
nonzero section and $bL-D$ is pseudoeffective for some integer $b>0$.
Choose finitely many admissible systems whose joint image has maximum
dimension $d$.  Every further admissible ratio is algebraic over this
image field.  Take its relative algebraic closure and a smooth
projective connected-fiber model $\pi:Z\to X$, $f:Z\to Y$, prepared
by flattening as before.  For every ample $H$ on $Y$ there is a $b_0>0$
with
\begin{equation}\label{conv:all-polarization}
 b_0\pi^*L-f^*H\quad\hbox{pseudoeffective}.
\end{equation}
Indeed the sum of the chosen divisors dominates the pullback of the
product polarization; its restriction to the joint image is big.
The sum itself is pseudoeffectively dominated by a multiple of $L$,
and a large multiple of the product polarization dominates $H$.
If $d=\dim X$, this makes $L$ big.

Take sections $u_j$ defining $N_j\sim m_jL-P$, put
$r=m_2-m_1$, and let $t=\pi^*(u_2/u_1)$, a rational section of
$r\pi^*L$.  For $j>2$ the expressions
$u_j^r u_1^{m_j-m_2}$ and $u_2^{m_j-m_1}$ belong to the same
admissible system, namely a multiple of $m_2L-P$.
Writing $D_t=\operatorname{div}(t)$ gives
\[
 r\pi^*N_j=r\pi^*N_1+(m_j-m_1)D_t+
                    f^*\operatorname{div}(g_j),
 \qquad g_j\in\mathbb C(Y)^*.
\]
Thus $D_t$ has nonnegative horizontal part, and its divisor of vertical
minima $M$ is pseudoeffective by the same fixed-error inequality used
in \eqref{conv:affine-divisors}.  The case $d=0$ is already settled.

For $d>0$, the full adjoint sheaf $f_*\mathcal O_Z(E)$,
$E=K_{Z/X}$, has rank one.  Otherwise two generic-fiber independent
sections can be made global in $E+kf^*H$.  After pushforward their
effective divisors give a subsystem of $k\pi_*f^*H$ on $X$ whose
ratio is not in $\mathbb C(Y)$.  By
\eqref{conv:all-polarization} this system is admissible, contradicting
maximality.  Mixing a positive-current metric furnished by
\eqref{conv:all-polarization} with the smooth pullback metric again
gives a globally defined coefficient $h$ with
$\Theta_h\ge\delta f^*\omega_Y$ and
$\mathcal J(h)=\mathcal O_Z$.  The singular direct-image theorem
therefore applies to the full actual sheaf as in the first proof.
Let $w$ be the pushed-forward Jacobian volume density on its smooth
base-change locus.  Then $-\log w$ has curvature at least
$\delta\omega_Y$ there, and $w$ has finite local mass.

Here the boundary can be written directly.  Set
\[
 c_Q=\min_{D\mapsto Q}
       \left\lfloor\frac{\operatorname{ord}_D E}
                            {\operatorname{ord}_D f^*Q}\right\rfloor,
 \qquad C=\sum_Qc_QQ.
\]
It is effective and integral.  If a prime over $Q$ is nonexceptional
then its coefficient in $E$ is zero, so $c_Q=0$.  Hence $f^*C$ is
exceptional, including the primes over codimension two by the model
condition.  Near a general point of a prime $D$ attaining the minimum,
write $f=(z_1^e,z_2,\ldots,z_d)$ and
$a=\operatorname{ord}_D E$.  A plurisubharmonic coefficient weight is
bounded above on a relatively compact chart.  Its metric norm is
therefore bounded below by a positive constant times a smooth norm.
Integration of this chart gives
\[
 w(y)\ge c|y_1|^{2(a+1-e)/e}.
\]
If $a=qe+s$, $0\le s<e$, then
$(a+1-e)/e=q+(s+1-e)/e\le q=\lfloor a/e\rfloor$.
It follows that
$\psi=-\log w+\log|s_C|^2$ is locally bounded above at general
points of every missing divisor.  Extend plurisubharmonically across
these divisors and then across codimension two, subtracting local
potentials of $\delta\omega_Y$ to preserve the lower bound.
This gives a metric on $-K_Y+C$ with global strict curvature and
\[
 |s_C|^2e^{-\psi}=w\in L^1_{\mathrm{loc}}.
\]
Theorem~\ref{thm:controlled-boundary} applies to $M$.
Lemma~\ref{conv:return}, now with the rational section $t$, completes
the alternative construction.  Thus its distinct ingredients are the
field of all admissible systems, full rank one from maximality, and the
explicit floor correction; the algebraic return is unchanged.

\subsection{From twisted tensors to a fixed error}
\label{conv:tensors}

The determinant method of Lazi\'c--Peternell
\cite[Lemma~4.1]{LP2018}, and its formulation in
\cite[Lemma~5.1]{LMPTX2023}, provides the link to twisted differentials.
Those cited extraction statements include a numerical-nontriviality
hypothesis.  The following direct construction also covers a numerically
trivial $L$, because the starting exponents here are positive and
unbounded.

\begin{lemma}\label{conv:determinant}
Let $X$ be smooth connected projective, let $L=-K_X$ be nef, and fix
$p>0$.  If
$H^0(X,(\Omega_X^1)^{\otimes p}\otimes mL)\ne0$ for unbounded
positive integers $m$, there are a pseudoeffective Cartier divisor $P$
and effective integral $N_j\sim m_jL-P$ for unbounded positive $m_j$.
If the starting sections are invariant for a torus and the induced
linearizations, the determinant sections defining $N_j$ are invariant.
\end{lemma}
\begin{proof}
Each section defines a map $\mathcal O_X(-mL)\to
(\Omega_X^1)^{\otimes p}$.  Saturate their generic span, obtaining a
coherent subsheaf $\mathcal S$.  Finitely many maps already span at
the generic point, and saturation ensures every original map factors
through $\mathcal S$.  Let $a=\operatorname{rank}\mathcal S$.
The reflexive determinant $A=(\bigwedge^a\mathcal S)^{**}$ is a line
bundle on smooth $X$.
The cotangent-subsheaf consequence of generic tangent nefness
\cite[Theorem~4.1]{LMPTX2023} gives $-A$ pseudoeffective.
More explicitly, saturate $A$ in the corresponding exterior power of
the tensor bundle and embed that exterior power in a cotangent tensor
power.  The negative of the saturated line is pseudoeffective; passing
back to $-A$ adds an effective divisor.  Put $P=-A$.

Choose a generic basis from finitely many starting sections.  Every
nonzero section of arbitrarily large twist can replace one vector of
that basis while keeping a basis.  On an unbounded subsequence the same
slot works.  Wedge it with the remaining fixed vectors.  This gives
nonzero global sections of $A\otimes m_jL$, where $m_j$ is the
varying exponent plus a fixed sum of exponents.  The sums remain
positive and unbounded.  The wedges extend globally by the determinant
map and reflexive extension across codimension two.
If the maps are equivariant, their saturation and determinant carry
the induced action, and the wedge of invariant sections is invariant.
\end{proof}

\begin{corollary}\label{conv:tensor-conversion}
Under the metric hypotheses of Theorem~\ref{conv:fixed-error}, if
$H^0(X,(\Omega_X^1)^{\otimes p}\otimes mL)\ne0$ for a fixed
$p\ge0$ and unbounded positive $m$, then a positive multiple of $L$
has a section.  The same holds with $\Omega_X^p$ in place of the
tensor power.
\end{corollary}
\begin{proof}
For $p=0$ a section is already given.  Otherwise apply
Lemma~\ref{conv:determinant} and Theorem~\ref{conv:fixed-error}.
Exterior forms embed in tensor powers by alternation in characteristic
zero.  This conversion is an ordinary statement; invariant determinant
sections alone do not turn its conclusion into an invariant section.
\end{proof}

\subsection{A movable-slope proof of the determinant input}
\label{conv:determinant-slope}

For completeness we give a second proof of the pseudoeffectivity used
in Lemma~\ref{conv:determinant}, under its smooth-metric application
hypotheses.  This proof explains the ramification correction in the
relative-volume argument.  It uses movable slope theory
\cite{GrebKebekusPeternell2016}, duality between pseudoeffective divisors
and movable curves \cite{BDPP2013}, and the positive-minimal-slope
algebraicity theorem of Campana--P\u aun
\cite[Theorem~1.1]{CampanaPaun2019}.

For a torsion-free sheaf $\mathcal E$ and movable curve class $\alpha$,
its slope is $\mu_\alpha(\mathcal E)=c_1(\mathcal E)\cdot\alpha/
\operatorname{rank}\mathcal E$.  The symbols $\mu_{\max,\alpha}$ and
$\mu_{\min,\alpha}$ denote the greatest and least slopes in its
Harder--Narasimhan filtration, whose successive quotients are
semistable with decreasing slopes.

Suppose a saturated subsheaf of a cotangent tensor power on smooth
projective $X$ had determinant of positive degree against a movable
curve class $\alpha$.  Tensor-product slope inequalities imply
$\mu_{\max,\alpha}(\Omega_X^1)>0$, equivalently
$\mu_{\min,\alpha}(T_X)<0$.  Since $c_1(T_X)\cdot\alpha\ge0$,
the positive-slope part $\mathcal F$ of the Harder--Narasimhan
filtration of $T_X$ satisfies
\[
 \mu_{\min,\alpha}(\mathcal F)>0,
 \qquad\mu_{\max,\alpha}(T_X/\mathcal F)\le0,
 \qquad c_1(T_X/\mathcal F)\cdot\alpha<0.
\]
The bracket obstruction from $\bigwedge^2\mathcal F$ to
$T_X/\mathcal F$ vanishes by the slope inequality.  Thus $\mathcal F$
is a foliation, and the cited algebraicity theorem realizes it as the
relative tangent foliation of a rational map with connected general
fibers.  Prepare an equidimensional model of that map, then take a
smooth resolution with maps $\mu:W\to X$ and $f:W\to Y$.
Put $L=-K_X$ and retain its smooth semipositive metric.

The relative adjoint bundle $K_{W/Y}+\mu^*L$ has a nonzero section on
a general fiber: there it is represented by the effective Jacobian
$K_{W/X}$.  On the smooth bundle and base-change locus, evaluation
with the $L^2$ direct-image metric supplies its Bergman weight
\[
 \Phi(x)=\log\sup_{\|u\|_{L^2}\le1}|u(x)|^2
\]
in a local relative-adjoint frame.  Semipositivity follows from the
dual-norm form of direct-image positivity
\cite{Berndtsson2009,Berndtsson2011}; zeros of the evaluation map give
allowed singularities.  On the isomorphism locus of $\mu$, the
submersive differential identifies this line with
$\det(T_X/\mathcal F)$.  The only extra codimension-one issue is
ramification.

A prime of $X$ vertical for the map dominates a base prime on the
normal equidimensional model.  Near its general point, and hence on
its strict transform in $W$, choose coordinates
\[
 f=(z_1^e,z_2,\ldots,z_{\dim Y}).
\]
The differential contributes $z_1^{e-1}$.  Therefore
$\det(T_X/\mathcal F)$ corresponds to the relative adjoint line
twisted down by $(e-1)\{z_1=0\}$.
A relative form written in total-space/base canonical frames acquires
the factor $(ez_1^{e-1})^{-1}$ when restricted to a fiber.
Integration on a fixed local fiber chart and the holomorphic submean
estimate consequently bound every unit-norm section coefficient by
$C|z_1|^{e-1}$.  It follows that
\[
 \Phi\le(e-1)\log|z_1|^2+O(1).
\]
After the indicated twist, the induced weight is locally bounded
above at the general point of every such prime.  At horizontal primes
the ordinary submean estimate gives the same upper boundedness without
a ramification term.  These weights extend plurisubharmonically across
codimension one, and then across the remaining codimension-two locus
on $X$.  They define a semipositive singular metric on
$\det(T_X/\mathcal F)$, making that line pseudoeffective.  This
contradicts its negative degree against $\alpha$.

By movable-cone duality, the negative determinant of every cotangent
tensor subsheaf is therefore pseudoeffective.  The argument retains the
factor $e-1$ because omitting it would construct a metric on the wrong
line at ramified divisors.

\section{Smooth coarea metrics and a dominated base}
\label{coa:analytic-tools}

Integration along a fibration turns an anticanonical metric upstairs into
an adjoint metric on its base.  In this section the coefficient metric is
smooth.  This regularity gives a direct curvature calculation and, by
coarea, an integrability exponent strictly greater than one.  Together
these facts produce a boundary correction supported only on base divisors
whose inverse images are exceptional over the original variety.

We first establish the analytic statement, keeping curvature and
integrability separate.  We then construct a maximal base dominated by
the anticanonical class and verify the metric hypotheses on that base.
The controlled-boundary criterion converts the resulting numerical
pseudoeffectivity into rational linear effectivity.

\subsection{The primitive representative and horizontal strictness}

\begin{lemma}\label{coa:curvature}
Let $f:\mathcal Z\to\mathbb D$ be a proper holomorphic submersion from a
K\"ahler manifold, with fibers of complex dimension $r>0$.  Let
$(\mathcal L,h)$ be a holomorphic line bundle with a smooth Hermitian metric
of semipositive curvature.  Suppose
$\mathcal E=f_*(K_{\mathcal Z/\mathbb D}+\mathcal L)$ is locally free and
commutes with base change.  Give $\mathcal E$ its fiberwise $L^2$ metric.
Write $D'$ for the $(1,0)$ part of the Chern connection of this
direct-image metric.  For a local holomorphic section $v$ with
$D'v(0)=0$, there is a smooth
absolute $\mathcal L$-valued $(r,0)$-form $w$ representing $v$ such that
\begin{equation}\label{coa:curvature-bound}
 -i\partial\bar\partial\|v\|^2\big|_0
 \ \geq\
 f_*\bigl(c_r w\wedge\overline w\,h\wedge i\Theta_{\mathcal L,h}\bigr)
 \big|_0,\qquad c_r=i^{r^2}.
\end{equation}
In particular, if $i\Theta_{\mathcal L,h}\geq c f^*(i\,dt\wedge d\bar t)$
on a nonempty open subset of the central fiber and its neighborhood,
where $c>0$ and $v(0)$ does not vanish identically on that subset, then the
curvature pairing of $\mathcal E$ with $v(0)$ is strictly positive at $0$.
\end{lemma}

\begin{proof}
We use the representative construction in Berndtsson's proof of direct-image
positivity \cite[Proposition~4.2 and formula~(4.8)]{Berndtsson2009}; the
one-dimensional-base formulation is also
\cite[formula~(2.1) and Lemma~2.1]{Berndtsson2011}.  The calculation below
makes clear why no strict positivity along the fibers is needed.

Write $F=f^{-1}(0)$, and fix a K\"ahler form $\omega$ on $\mathcal Z$.
A smooth splitting of the relative cotangent sequence gives an absolute
lift $w$ of the relative form $v$.  Since $dt\wedge w$ is holomorphic, write
\[
 \bar\partial w=dt\wedge\eta,
 \qquad \partial^h w=dt\wedge\mu,
\]
where $\partial^h$ is the $(1,0)$ part of the Chern connection.  On $F$,
$\eta$ has type $(r-1,1)$ and $\mu$ has type $(r,0)$.
Differentiating the fiberwise pairings of $v$ with local holomorphic
sections of $\mathcal E$ identifies the projection of $\mu$ onto the
holomorphic $\mathcal L|_F$-valued top forms with $D'v(0)$.  These sections
exhaust the top forms on $F$ by base change.  Thus $\mu|_F$ is orthogonal
to the holomorphic top forms.  The terms obtained by differentiating the
conjugate relative section have a $d\bar t$ factor and do not contribute
to this $(1,0)$ derivative.

Furthermore, $(\omega\wedge\eta)|_F$ is $\bar\partial$-exact.  Indeed
$w\wedge\omega$, of type $(r+1,1)$ on the total space, is divisible by
$dt$; applying $\bar\partial$ and using $d\omega=0$ gives the assertion
on $F$.  We can therefore change the lift to $w+dt\wedge\beta$, where
$\beta$ has type $(r-1,0)$, so that
\begin{equation}\label{coa:primitive-choice}
 \mu|_F=0,\qquad (\omega\wedge\eta)|_F=0
\end{equation}
for the changed lift.  To verify simultaneous solvability, put
$\chi=\omega\wedge\beta$.  Wedge multiplication by $\omega$ is an
isomorphism from $(r-1,0)$-forms to $(r,1)$-forms on $F$, and the
K\"ahler identity gives
\[
 \bar\partial^*\chi=i\partial^h\beta.
\]
The required equations are
$\bar\partial^*\chi=i\mu$ and
$\bar\partial\chi=\omega\wedge\eta$.
The first right-hand side lies in the image of $\bar\partial^*$ because
of the preceding orthogonality; the second is $\bar\partial$-exact.
Let $G_{\bar\partial}$ denote the Green operator of the coefficient
$\bar\partial$-Laplacian on the compact fiber.  With
$q=\omega\wedge\eta$, the two forms
\[
 \chi_1=\bar\partial G_{\bar\partial}(i\mu),\qquad
 \chi_2=\bar\partial^*G_{\bar\partial}q
\]
satisfy
$\bar\partial^*\chi_1=i\mu$, $\bar\partial\chi_1=0$,
$\bar\partial\chi_2=q$, and $\bar\partial^*\chi_2=0$.
Indeed $i\mu$ is orthogonal to the harmonic $(r,0)$-forms and $q$ is
$\bar\partial$-exact.  Thus $\chi=\chi_1+\chi_2$ solves both equations.
Inverting $\omega\wedge$ and extending $\beta$ smoothly off $F$ proves
\eqref{coa:primitive-choice}.  When $r=1$, $q=0$ by type and the same
construction applies.

For this lift, differentiation of
$\|v\|^2=c_r f_*(w\wedge\overline w\,h)$ gives
\[
 \partial\|v\|^2
   =c_r f_*(\partial^h w\wedge\overline w\,h).
\]
Apply $i\bar\partial$, use
$\bar\partial\partial^h w
 =\Theta_{\mathcal L,h}\wedge w-\partial^h\bar\partial w$,
and note that the quadratic term in $\partial^h w$ vanishes on $F$.
To handle the remaining derivative, apply $\partial$ to
$f_*(\bar\partial w\wedge\overline w\,h)=0$; that pushforward vanishes
by type and its $dt$ factor.  At $0$ the resulting identity is
\begin{equation}\label{coa:primitive-identity}
 \begin{split}
 i\bar\partial\partial\|v\|^2
 &= f_*\bigl(c_r w\wedge\overline w\,h
                         \wedge i\Theta_{\mathcal L,h}\bigr)\\
 &\quad-i\,dt\wedge d\bar t\,
                  c_r\int_F\eta\wedge\overline\eta\,h.
 \end{split}
\end{equation}
Metric pairing is denoted multiplicatively by $h$ in these formulas.
For a primitive $(r-1,1)$-form the last term has nonnegative sign and is
a positive normalization of its squared norm.  This proves
\eqref{coa:curvature-bound}.

The integrand on the right of \eqref{coa:curvature-bound} is nonnegative:
a form of type $(r,0)$ in complex dimension $r+1$ is decomposable, and
wedging its squared form with a semipositive $(1,1)$-form gives a
nonnegative top-degree form.  On the stated patch we can replace the
curvature by $c f^*(i\,dt\wedge d\bar t)$ in a lower bound.  Only the
relative restriction of $w$ then contributes.  Its integral over an open
set on which $v(0)$ is nonzero is positive.  This proves strictness.
\end{proof}

For a higher-dimensional base we apply Lemma~\ref{coa:curvature} to small
embedded disks in its tangent directions.  In rank one a local holomorphic
frame may be chosen Chern-normal at a specified point.  Dividing
\eqref{coa:curvature-bound} by its squared norm then gives the curvature
of minus the logarithm of that norm.  A lower bound by the pullback of a
positive base form on one patch yields strict positivity in every
nonzero tangent direction at the image point.  If the relative dimension
is zero and the generic rank is one, the smooth map is locally an
isomorphism and this assertion follows directly from the coefficient
metric.

\subsection{A coarea estimate}

\begin{lemma}\label{coa:coarea}
Let $f:Z\to Y$ be a surjective morphism between connected smooth projective
varieties, with $\dim Z=n$ and $\dim Y=b>0$.  Let $\mu$ be a smooth
nonnegative top-degree form on $Z$.  Then $f_*\mu$ is absolutely continuous
with respect to any smooth positive volume on $Y$, and its density belongs
to $L^{1+\eta}(Y)$ for some $\eta>0$.
\end{lemma}

\begin{proof}
Fix K\"ahler metrics on $Z$ and $Y$, and use their volumes and the induced
fiber volumes.  On the smooth locus of $f$ let $J_f$ be its real coarea
Jacobian.  In orthonormal complex frames $J_f$ is the sum of the squared
absolute values of the top complex minors of $df$; equivalently, up to a
fixed normalization,
\[
 f^*\omega_Y^b\wedge\omega_Z^{n-b}=J_f\,\omega_Z^n.
\]
In finitely many coordinate charts it is comparable with the squared
norm of the corresponding holomorphic Jacobian-minor ideal.  This ideal
is not identically zero, because $f$ is dominant in characteristic zero.
A log resolution of the ideal reduces small inverse powers to products
of powers of coordinate absolute values.  Choosing their exponents
smaller than their finitely many integrability thresholds, and using
compactness, gives
\begin{equation}\label{coa:inverse-jacobian}
 \int_Z J_f^{-\eta}\,dV_Z<\infty
\end{equation}
for some $\eta>0$.

Let $U\subset Y$ be a dense Zariski open over which $f$ is smooth.  Since
$\mu\leq C\,dV_Z$, the coarea formula gives, for the density $\rho$,
\[
 \rho(y)\leq C\int_{Z_y}J_f^{-1}\,dV_{Z_y}
 \qquad (y\in U).
\]
The fiber volumes are constant on each connected smooth-base component,
by closedness of $\omega_Z^{n-b}$ and proper smooth local triviality.
They are in particular uniformly bounded.  The power-mean inequality
and coarea therefore give
\begin{equation}\label{coa:density-estimate}
 \begin{split}
 \int_U\rho^{1+\eta}\,dV_Y
 &\leq C'\int_U\int_{Z_y}J_f^{-1-\eta}\,dV_{Z_y}\,dV_Y\\
 &=C'\int_{f^{-1}(U)}J_f^{-\eta}\,dV_Z<\infty.
 \end{split}
\end{equation}
Finally, $f^{-1}(Y\setminus U)$ is a proper analytic subset of $Z$, so
has zero measure for the smooth form $\mu$.  There is consequently no
additional measure supported on $Y\setminus U$.  This proves both
absolute continuity and the claimed exponent.
\end{proof}

\subsection{Extension across base divisors}

\begin{proposition}\label{coa:volume-metric}
Let $X$ and $Z$ be connected smooth projective varieties, let
$g:Z\to X$ be a birational morphism, and let $f:Z\to Y$ be a surjective
morphism to a connected smooth projective variety of positive dimension.
Put $L=-K_X$, and assume $L$ has a smooth Hermitian metric $h$ with
semipositive curvature.  Let $E=K_{Z/X}$ be the effective Jacobian divisor.
Assume
\[
 \operatorname{rank} f_*\mathcal O_Z(E)=1.
\]
Let $\mathcal S$ be the set of prime divisors $G$ on $Y$ such that every
prime divisor of $Z$ dominating $G$ is $g$-exceptional.  Suppose there is
an open patch over the smooth base-change locus of $f$, disjoint from
$E$, on which the curvature of $g^*h$ dominates a positive multiple of
the pullback of a K\"ahler form on $Y$.
Then there is an effective integral divisor $A$ supported on
$\mathcal S$ and a singular metric on $M=-K_Y+A$, with psh weights $\psi$,
such that the metric is smooth and strictly positively curved on a
nonempty open subset of $Y$ and
\begin{equation}\label{coa:weighted-lp}
 e^{-\psi}|s_A|^2\in L^{1+\eta}_{\mathrm{loc}}
 \qquad\text{for some }\eta>0.
\end{equation}
Here $|s_A|$ means the absolute value of the coefficient of the canonical
section in a local holomorphic frame.
\end{proposition}

\begin{proof}
The identity section of $K_X+L$ and $h$ determine a smooth positive volume
form $\mu_X$ on $X$.  Write $\rho$ for the local density of
$f_*(g^*\mu_X)$ in base coordinates.  On a dense open $U\subset Y$, the
map $f$ is smooth, base change holds, and the canonical section $s_E$
spans $H^0(Z_y,\mathcal O_{Z_y}(E))$.  It gives the identification
\[
 f_*(K_{Z/Y}+g^*L)|_U\simeq -K_Y|_U.
\]
The fiberwise $L^2$ norm of the corresponding coordinate anticanonical
frame is $\rho$.  This follows directly by contracting $s_E$ with its
conjugate through $g^*h$ and integrating along the fibers; the resulting
top-degree form is $g^*\mu_X$.  Lemma~\ref{coa:curvature} proves that
$-\log\rho$ is psh, and the assumed patch proves strictness at its image.
Smooth dependence of fiber integrals gives smoothness on $U$.

It remains to control the metric at the boundary.  Let $G$ be a prime
component of $Y\setminus U$, and choose a prime $J\subset Z$ dominating
$G$.  Put $e_J=\operatorname{mult}_J(f^*G)>0$ and
$h_J=\operatorname{mult}_J E$.  Near general points of $J$ and $G$ there
are coordinates with
\[
 G=(t_1=0),\quad J=(z_1=0),\qquad
 t_1=z_1^{e_J},\quad t_i=z_i\ (2\leq i\leq b).
\]
Indeed $J\to G$ is generically submersive, and the unit in the first
coordinate can be absorbed into $z_1$.  The remaining coordinates are
fiber coordinates.  On such a patch, the pullback volume is bounded below
by a positive constant times $|z_1|^{2h_J}$ times coordinate volume.
Changing variables in the first coordinate, and integrating over a
fixed small fiber-coordinate polydisk, gives
\begin{equation}\label{coa:ramification-lower}
 \rho(t)\geq c\,|t_1|^{2((h_J+1)/e_J-1)}
\end{equation}
near general points of $G$, on $U$.  The denominator here is the squared
Jacobian $e_J^2|z_1|^{2(e_J-1)}$.

If $G\notin\mathcal S$, choose $J$ nonexceptional for $g$.  Then $h_J=0$,
and the exponent in \eqref{coa:ramification-lower} is nonpositive.
Thus $-\log\rho$ is locally bounded above there.  If
$G\in\mathcal S$, choose an integer
$a_G\geq\max\{0,(h_J+1)/e_J-1\}$.  With
$A=\sum_{G\in\mathcal S}a_GG$, the weight
\[
 \psi=-\log\rho+\log|s_A|^2
\]
is bounded above at general boundary points.  Only finitely many
$G\in\mathcal S$ occur, since each is the image of one of the finitely
many $g$-exceptional primes.  On $U$ the added term is pluriharmonic.
The psh extension theorem therefore extends the weight across all
codimension-one points under consideration.  The remaining analytic
subset has codimension at least two, across which psh extension is
again available.  One can see the needed local upper bound by using
small disks through nearby points whose boundaries remain in a fixed
compact subset off that analytic set, and applying the maximum
principle.  The frame transformations identify these extended weights
as a metric of $-K_Y+A$.

The strict patch can be chosen away from $A$, so strict positivity is
preserved there.  Lastly, $g^*\mu_X$ is a smooth nonnegative volume form
on the compact variety $Z$.  Lemma~\ref{coa:coarea} gives an exponent
$1+\eta>1$ for its pushforward density.  Since
$e^{-\psi}|s_A|^2=\rho$ almost everywhere in the compatible frames,
this is exactly \eqref{coa:weighted-lp}.
\end{proof}

\begin{corollary}[Smooth invariant direct image]\label{coa:invariant-volume-metric}
The conclusion of Proposition~\ref{coa:volume-metric} remains valid if
its ordinary rank-one hypothesis is replaced by
\[
 \operatorname{rank}\bigl(f_*\mathcal O_Z(K_{Z/X})\bigr)^T=1,
\]
provided that a compact torus $T$ acts holomorphically on $X$ and $Z$,
acts trivially on $Y$, preserves $f$, $g$, and the smooth metric $h$,
and all canonical bundles carry their natural linearizations.
\end{corollary}

\begin{proof}
On the smooth base-change locus, averaging over normalized Haar measure
is a holomorphic idempotent on the adjoint direct image.  Invariance of
the fiber metric makes this projection orthogonal.  Its image is the
invariant line, and its kernel is the orthogonal holomorphic complement.
Consequently the line has the quotient metric of the semipositive
smooth direct image.  For a section in this line, a Chern-normal
holomorphic extension within the line is also Chern-normal in the full
bundle, so Lemma~\ref{coa:curvature} proves the same strictness statement.
The naturally linearized Jacobian section is invariant and generates
this line.  Its norm is still the density $\rho$.

For relative dimension zero the full direct image is locally a direct
sum of the coefficient lines on the sheets of an unramified cover.
Its curvature is their direct sum; the same orthogonal projection
argument applies, and the nonzero Jacobian section detects the strict
patch.  The ramification estimate, divisor correction, and coarea
calculation in Proposition~\ref{coa:volume-metric} do not use ordinary
rank one.  They therefore prove the stated extension and integrability.
\end{proof}

\subsection{From local strictness to a global strict metric}

The metric just constructed is strictly positive on one open set.
The next step uses the extra integrability exponent to introduce a
global lower curvature bound while preserving the distinguished
integrable section.

\begin{lemma}\label{coa:weighted-effectivity}
Let $Y$ be smooth connected projective of positive dimension, let
$A\geq0$ be integral, and put $M=-K_Y+A$.  Suppose $M$ has psh metric
weights $\psi$, smooth and strictly positively curved on a nonempty
open set, and
\[
 |s_A|^2e^{-\psi}\in L^{1+\eta}_{\mathrm{loc}}
 \quad\text{for some }\eta>0.
\]
Then $M$ has a singular metric with a global strict curvature lower
bound and with $|s_A|^2$ times its squared frame norm locally
integrable.  For every pseudoeffective rational divisor $W$, an
effective rational divisor $C_A$ supported on $A$ makes $W+C_A$
rationally linearly equivalent to an effective rational divisor.
\end{lemma}

\begin{proof}
The curvature current of $\psi$ has positive absolutely continuous
top-degree mass, since it is smooth and positive definite on the
specified open set.  Boucksom's volume criterion
\cite{Boucksom2002} makes $M$ big.  Write
$M\sim_{\mathbb Q}H+G$ with $H$ ample rational and $G\geq0$ rational.
A smooth positive metric on $H$ and the divisor metric of $G$ give a
metric on $M$ with weights $\chi$, curvature bounded below by a
positive multiple of a K\"ahler form, and divisorial analytic
singularities.  Log resolution and compactness give a number
$\tau>0$ such that $|s_A|^2e^{-\chi}\in L^\tau_{\mathrm{loc}}$.
Choose $\sigma\in(0,1)$ sufficiently small that
\[
 \frac{1-\sigma}{1+\eta}+\frac{\sigma}{\tau}<1.
\]
The weights $\psi_\sigma=(1-\sigma)\psi+\sigma\chi$ then have a
global strict curvature lower bound.  Moreover,
\[
 |s_A|^2e^{-\psi_\sigma}
 =\bigl(|s_A|^2e^{-\psi}\bigr)^{1-\sigma}
  \bigl(|s_A|^2e^{-\chi}\bigr)^\sigma
\]
is locally integrable by H\"older's inequality.  Equivalently,
$\mathcal J(\psi_\sigma)\supset\mathcal O_Y(-A)$.
The controlled-boundary criterion, Theorem~\ref{thm:controlled-boundary},
now gives the assertion about $W$ and $C_A$.
\end{proof}

\subsection{The smooth metric on the maximal dominated base}
\label{coa:maximal-base}\label{coa:dominated-conversion}

We now apply the smooth coarea construction to the maximal base of
Section~\ref{conv:all-dominated}.  That base already supplies the
interpolation and the pseudoeffective divisor of vertical minima.
The additional points here are a valuation proof of the model
condition and a direct intersection argument for horizontal strictness.
These allow us to use the original smooth metric, and hence retain
the exponent $1+\eta$ in Proposition~\ref{coa:volume-metric}.

Retain the hypotheses of Theorem~\ref{conv:fixed-error}, with effective
integral divisors $N_j\sim a_jL-D$, where $D$ is pseudoeffective and
$a_j$ is a strictly increasing unbounded sequence of positive integers.
A rational map is \emph{dominated by $L$} if, on a smooth resolution
$X\xleftarrow{\ g\ }Z\xrightarrow{\ f\ }Y$, one has
\[
 g^*L-\epsilon f^*H\quad\hbox{pseudoeffective}
\]
for an ample $H$ on the projective base and a rational $\epsilon>0$.
Include the point map.  This is the map formulation of the admissible
systems used in Section~\ref{conv:all-dominated}.  Indeed, an
admissible system gives this inequality after resolving its base
ideal.  Conversely, after multiplying $H$ so that it defines a
projective embedding, its pulled-back sections push down to a
subsystem of $|g_*f^*H|$; pushing down the domination inequality
makes that system admissible.

The condition persists on higher source resolutions and under
birational base modifications.  It also persists under a generically
finite change of base to which the rational map lifts dominantly:
the pullback of the old ample class is big and dominates a small
ample class.  Finally, the joint image of two dominated maps is
dominated.  On a common source resolution, add suitable positive
multiples of their inequalities and use the sum of the pulled-back
polarizations, whose restriction to the joint image is ample.
Thus the maximality construction in Section~\ref{conv:all-dominated}
provides a base field $\mathbb C(Y)$, relatively algebraically closed
in $\mathbb C(X)$, containing the functions supplied by every
dominated map.

For the resulting smooth projective diagram put
\[
 P=\frac{g^*N_2-g^*N_1}{a_2-a_1}\sim_{\mathbb Q}g^*L,
 \qquad B=a_1P-g^*N_1.
\]
The interpolation argument of Section~\ref{conv:all-dominated} gives
$\mathbb Q$-principal divisors $U_j$ on $Y$ with
\begin{equation}\label{coa:interpolation}
 g^*N_j=a_jP-B+f^*U_j,\qquad U_1=U_2=0.
\end{equation}
For $j>2$, the two effective divisors in
$(a_j-a_1)N_2\sim(a_j-a_2)N_1+(a_2-a_1)N_j$ lie in a
multiple of the admissible system $N_2$, so their ratio belongs to
$\mathbb C(Y)$.  In particular $P$ has nonnegative horizontal
part.  If $Y$ is a point, pushing $P$ down already proves the theorem.
Assume henceforth that $b=\dim Y>0$.

\subsubsection*{Preparing vertical primes by their valuations}

The model can be arranged so that every vertical prime not exceptional
over $X$ dominates a base prime.  The flattening argument of
Section~\ref{conv:section} gives one construction.  The following
valuation argument gives another and explains why a base modification
suffices.

On an initial model only finitely many source primes map into base
codimension at least two: take an open set over which $f$ is
equidimensional and inspect the divisorial components of its
complement.  The valuation of such a prime restricts nontrivially to
$\mathbb C(Y)$, since a local function vanishing on its proper center
has positive order.  The restricted discrete valuation is divisorial.
To see this, its residue transcendence degree can drop by at most
$\dim X-\dim Y$.  Residues algebraically independent over the
restricted residue field lift to elements algebraically independent
over $\mathbb C(Y)$: a polynomial relation, with coefficients scaled
to have minimum valuation zero, would give a residue relation.
The source residue field has transcendence degree $\dim X-1$.
The restricted residue field therefore has transcendence degree at
least $b-1$, and the valuation inequality gives the opposite bound.

Choose a projective base model on which these finitely many valuations
have divisorial centers, and resolve it.  One can construct such a
model by resolving maps to copies of $\mathbb P^1$ defined by
$b-1$ independent residues: the center then has dimension at least
$b-1$ and remains proper.  Take a common smooth resolution of the
previous source and the graph over the new base.  Primes already
dominating base divisors retain that property.  New source primes
exceptional over the old source are also exceptional over $X$.
This proves the required preparation.  The domination and field
properties persist by the preceding observations.  Retain $P,B$ for
their pullbacks to this model and $U_j$ for the corresponding principal
divisors on the modified base; \eqref{coa:interpolation} is unchanged.

On this prepared model define
\begin{equation}\label{coa:minima}
 w_G=\min_{J\mapsto G}\frac{\operatorname{mult}_JP}
                              {\operatorname{mult}_Jf^*G},
 \qquad W=\sum_Gw_GG.
\end{equation}
The sum has finite support.  Formula~\eqref{coa:interpolation} makes
$W$ pseudoeffective by the same fixed-error limit as in
Section~\ref{conv:all-dominated}: if
$c_G=\min_{J\mapsto G}\operatorname{mult}_JB/
\operatorname{mult}_Jf^*G$ and $C=\sum_Gc_GG$, then
$a_jW+U_j-C\ge0$; dividing by $a_j$ gives numerical classes
converging to $[W]$.

Let $\mathcal S$ consist of base primes whose dominating source
primes are all $g$-exceptional.  It is finite, and the model condition
makes the entire pullback of each member exceptional over $X$.
Lemma~\ref{conv:return}, with its $R=P$ and $M=W$, reduces the
proof to making $W+C_{\mathcal S}$ rationally effective for an
effective correction supported on $\mathcal S$.

\subsubsection*{Rank one and horizontal strictness}

Put $E=K_{Z/X}$.  The maximality argument in
Section~\ref{conv:all-dominated} gives
\begin{equation}\label{coa:rank-one}
 \operatorname{rank}f_*\mathcal O_Z(E)=1.
\end{equation}
Its concrete input is a domination inequality
$b'g^*L-f^*H'$ pseudoeffective for a sufficiently ample base
polarization $H'$.  Two generic-fiber independent adjoint sections
would, after an ample base twist, push down to an admissible pencil
with ratio outside $\mathbb C(Y)$, contradicting the defining
maximality.  The Jacobian section supplies the nonzero rank-one
generator.

Let $\alpha\ge0$ be the smooth curvature of the original pulled-back
metric on $g^*L$, normalized to represent $g^*c_1(L)$, and let $r$
be its maximum rank.  If $r=n=\dim Z$, then $L$ is nef with positive
top self-intersection and hence big, giving the required nonvanishing.
Suppose $r<n$.  Choose K\"ahler
forms $\omega_Z,\omega_Y$ representing ample classes.
Domination gives $b'[\alpha]-[f^*\omega_Y]$ pseudoeffective for
some $b'>0$.  Intersecting with
$[\alpha]^r[\omega_Z]^{n-r-1}$, and using $\alpha^{r+1}=0$, yields
\[
 \int_Z f^*\omega_Y\wedge\alpha^r\wedge\omega_Z^{n-r-1}=0.
\]
The integrand is nonnegative, so it vanishes pointwise.  At a
rank-$r$ point, this says that $f^*\omega_Y$ vanishes on
$\ker\alpha$.  A nonnegative Hermitian form has zero mixed pairings
against any vector on which its quadratic value vanishes.  Hence
\begin{equation}\label{coa:kernel-inclusion}
 \ker\alpha\subseteq\ker df.
\end{equation}
At a smooth point this forces $r\ge b$; rank zero is therefore
excluded.  The maximum-rank open set meets both the complement of
$E$ and the smooth base-change locus.  On a smaller relatively compact
constant-rank patch there, \eqref{coa:kernel-inclusion} gives
$\alpha\ge c f^*\omega_Y$ for some $c>0$.

Proposition~\ref{coa:volume-metric} now applies with its birational
map $g:Z\to X$ and fibration $f:Z\to Y$.  It supplies an integral
effective $A$ supported on $\mathcal S$ and a semipositive metric
on $-K_Y+A$, smooth and strictly positive on a nonempty open set,
whose weighted density belongs to $L^{1+\eta}_{\rm loc}$.
Lemma~\ref{coa:weighted-effectivity} gives an effective $C_A$
supported on $A$ with $W+C_A$ rationally effective.
The return lemma applies to $P\sim_{\mathbb Q}g^*L$: its horizontal
part is nonnegative, its minimum divisor is $W$, and every component
of $f^*C_A$ is exceptional over $X$.  It gives the desired section
of a positive multiple of $L$.

This proves Theorem~\ref{conv:fixed-error} using the original smooth
coefficient metric.  The determinant and tensor consequences remain
Lemma~\ref{conv:determinant} and Corollary~\ref{conv:tensor-conversion}.
The next section retains a different bigness proof and a
fiber-intersection test for horizontal positivity.

\section{Jet separation and a nef-and-big decomposition}
\label{jet:iitaka-section}

There is another way to obtain the boundary needed for nonvanishing.
A metric that is strictly positive on one open set first separates
linearly growing jets at a point.  This proves bigness directly by Nadel
vanishing.  A subsequent multiplier-ideal construction produces a nef
and big divisor on a resolution, which can be perturbed by any fixed
pseudoeffective class.  We prove this base lemma independently, then
apply it to an Iitaka fibration.  The Iitaka construction also provides
a different proof of the horizontal curvature bound.

\subsection{A base lemma proved by jets and a nef-and-big decomposition}

\begin{theorem}\label{jet:base-lemma}
Let $Y$ be a smooth connected projective variety of dimension $b>0$,
let $A\geq0$ be an integral divisor, and put $M=-K_Y+A$.
Assume $M$ has a singular Hermitian metric with psh local weights
$\varphi$, smooth and strictly positively curved on a nonempty open
subset, and satisfying
\begin{equation}\label{jet:weighted-assumption}
 e^{-\varphi}|s_A|^2\in L^{1+\gamma}_{\mathrm{loc}}
 \qquad\text{for some }\gamma>0.
\end{equation}
For every pseudoeffective $\mathbb Q$-divisor $W$ on $Y$ there is an
effective $\mathbb Q$-divisor $C_A$ supported on $\operatorname{Supp}A$
such that $W+C_A$ is $\mathbb Q$-linearly equivalent to an effective
$\mathbb Q$-divisor.
\end{theorem}

\begin{proof}
We first prove that $M$ is big, using its strict positivity only on the
specified patch.  Choose ample Cartier divisors $H_1,A_1$ so that
$B_1+A_1$ is very ample, where $B_1=K_Y+H_1$.
At a point $y_0$ in the patch, use coordinates $z$ and a smooth cutoff
to form a function $\ell$ equal to $\log\|z\|^2$ near $y_0$, supported
in a ball relatively compact in that patch, and zero outside it.
Fix a background K\"ahler form $\omega$.  On the support of the
cutoff error choose constants $\delta>0$ and $C>0$ such that
$i\partial\bar\partial\varphi\geq\delta\omega$ and
$i\partial\bar\partial\ell\geq-C\omega$; outside that support the
added logarithmic pole has nonnegative curvature.  Fix
$0<c\leq\delta/(2C)$.  For large integers $h$ and
$N=\lfloor ch\rfloor$, the metric
on $hM+H_1$ with weight
\[
 h\varphi+N\ell+\varphi_{H_1}
\]
therefore has curvature bounded below by a K\"ahler form.
Indeed the strict curvature of $h\varphi$ absorbs the cutoff error,
and the positive metric on $H_1$ supplies a global lower bound.

Let $\mathcal I_h$ be this metric's multiplier ideal.  Nadel vanishing
\cite{Nadel1990} gives
\[
 H^1\bigl(Y,\mathcal O_Y(hM+B_1)\otimes\mathcal I_h\bigr)=0.
\]
Near $y_0$, the original weight is smooth and
\[
 \mathcal I_{h,y_0}=\mathfrak m_{y_0}^{N-b+1}
\]
for $N\geq b$.  In fact a holomorphic monomial of total order $q$
is integrable against $\|z\|^{-2N}$ exactly when $q>N-b$.
The support of the quotient by this ideal has $y_0$ as an isolated
component, separated from its other support by a neighborhood on which
the original weight is smooth.  Any jet in
$\mathcal O_{Y,y_0}/\mathfrak m_{y_0}^{N-b+1}$ thus defines a global
section of this quotient, taken to be zero on the other components.
The vanishing lifts every such jet.  The quotient has length
$\binom{N}{b}$, so in fact
\[
 h^0(Y,hM+B_1)\geq\binom{\lfloor ch\rfloor}{b}.
\]
Consequently
\begin{equation}\label{jet:section-growth}
 h^0(Y,hM+B_1)\geq c' h^b
\end{equation}
for a constant $c'>0$ and all sufficiently large $h$.

Take a fixed smooth divisor $T_0\in|B_1+A_1|$.  We have
$h^0(T_0,(hM+B_1)|_{T_0})=O(h^{b-1})$.
For completeness, choose an ample divisor $H_2$ such that
$H_2-M$ has a section not identically zero on any component of $T_0$;
its $h$-th power embeds these restricted sections into those of
$(hH_2+B_1)|_{T_0}$.  The latter dimensions have the asserted growth by
the ample Hilbert polynomial, including the bounded zero-dimensional
case when $b=1$.  The restriction exact sequence and
\eqref{jet:section-growth} now imply
\[
 H^0(Y,hM-A_1)\ne0
\]
for large $h$.  Thus $M$ is big.

Write $M$ as the sum of an ample rational class and an effective
rational divisor.  This gives another metric with weight $\psi$,
curvature bounded below by a positive multiple of a K\"ahler form,
and divisorial analytic singularities.  Small inverse powers of its
local divisor equation are integrable, so
$e^{-\psi}|s_A|^2\in L^\tau_{\mathrm{loc}}$ for some $\tau>0$.
Replace $\varphi$ by $(1-\sigma)\varphi+\sigma\psi$, for a sufficiently
small $\sigma>0$.  Its curvature now has a global positive lower bound.
It also retains the local integrability
\begin{equation}\label{jet:mixed-l1}
 e^{-\varphi}|s_A|^2\in L^1_{\mathrm{loc}}.
\end{equation}
Indeed the new expression is the product of the two old weighted
expressions to the powers $1-\sigma$ and $\sigma$, and H\"older applies
as soon as
$(1-\sigma)/(1+\gamma)+\sigma/\tau<1$.
We keep the notation $\varphi$ for this new weight.

We now extract a nef-and-big divisor on a resolution, while retaining
the integrability needed to bound every discrepancy.
Fix a very ample divisor $H$ with a positive metric.  Take an integer
$k>0$ sufficiently large that $P=kM-H$ has a semipositive metric with
weight $k\varphi-\varphi_H$.  Put
\[
 B_0=K_Y+(b+1)H,
 \qquad C=mP+B_0,
\]
where $m$ is a positive integer so large that $mH-B_0$ is ample.
Set $\mathcal J=\mathcal I(mk\varphi)$.
For $1\leq j\leq b$, Nadel vanishing applied to
$mP+(b+1-j)H$ gives
\[
 H^j\bigl(Y,\mathcal O_Y(C-jH)\otimes\mathcal J\bigr)=0.
\]
The added smooth weights leave the multiplier ideal unchanged, and
the curvature is strictly positive because $b+1-j\geq1$.
Castelnuovo--Mumford regularity therefore makes
$\mathcal O_Y(C)\otimes\mathcal J$ globally generated.

Resolve $\mathcal J$ and $A$ by a smooth projective birational morphism
$d:Z\to Y$, including the exceptional divisor in the resolved normal
crossings support.  Write
$\mathcal J\mathcal O_Z=\mathcal O_Z(-G)$.
Then $d^*C-G$ is globally generated, and
\begin{equation}\label{jet:nef-big-decomposition}
 d^*M\sim_{\mathbb Q}Q+\frac{G}{mk},\qquad
 Q=\frac{d^*C-G+d^*(mH-B_0)}{mk}.
\end{equation}
The divisor $Q$ is nef and big: its first summand is nef, and the
pullback of the ample divisor $mH-B_0$ is nef and big.
This decomposition will allow us to perturb by $d^*W$ while preserving
control of the boundary coefficients.

Put
\begin{equation}\label{jet:delta-zero}
 \Delta_0=\frac{G}{mk}-K_{Z/Y}-d^*A.
\end{equation}
We claim that $\Delta_0^+$ is a klt boundary.  At a general point of
a prime $R=(z_1=0)$ in the resolved support, let $r,a,c$ be its
coefficients in $G,K_{Z/Y},d^*A$, respectively.
The Ohsawa--Takegoshi point-extension theorem
\cite{OhsawaTakegoshi1987} gives
\begin{equation}\label{jet:point-estimate}
 e^{mk\varphi(d(z))}\leq C_1|z_1|^{2r}.
\end{equation}
Explicitly, point extension gives holomorphic functions with uniformly
bounded $mk\varphi$-weighted $L^2$ norm and
$|q_{y'}(y')|^2\geq c_1e^{mk\varphi(y')}$ for a fixed $c_1>0$.  Local upper bounds for the weight give
ordinary sup bounds.  Pullback makes every such function divisible by
$z_1^r$; uniform Cauchy bounds on the quotient and $y'=d(z)$ give
\eqref{jet:point-estimate}.  This argument uses the coherent multiplier
ideal itself, and requires no analytic-singularity hypothesis on
$\varphi$.

Changing variables in \eqref{jet:mixed-l1} includes the squared
Jacobian of $d$, of order $2a$.  Inequality
\eqref{jet:point-estimate} then forces integrability of
$|z_1|^{2(a+c-r/(mk))}$.  Hence
\begin{equation}\label{jet:strict-discrepancy}
 \frac r{mk}<1+a+c.
\end{equation}
All coefficients of $\Delta_0$ are consequently less than one.
Its support has simple normal crossings, so $(Z,\Delta_0^+)$ is klt.
Moreover \eqref{jet:nef-big-decomposition} and $M=-K_Y+A$ imply
\[
 K_Z+\Delta_0\sim_{\mathbb Q}-Q.
\]

It remains to incorporate $W$ without losing the strict klt inequalities.
Write $Q\sim_{\mathbb Q}A_0+G'$ with $A_0$ ample and $G'\geq0$.
For a small rational $\lambda\in(0,1)$,
\[
 Q-\lambda G'\sim_{\mathbb Q}(1-\lambda)Q+\lambda A_0
\]
is ample.  Check the pair on a common log resolution of $\Delta_0^+$ and
$G'$.  Its crepant boundary coefficients depend affinely on $\lambda$
and are strictly below one at $\lambda=0$.  Thus
$(Z,\Delta_0^++\lambda G')$ remains klt for sufficiently small
$\lambda$.  This resolution is used only to check discrepancies;
$Q$ and $\Delta_0$ continue to denote the divisors on $Z$.  By openness of the ample cone choose a further
small rational $\epsilon>0$ so that
$Q-\lambda G'+\epsilon d^*W$ is ample.
Take a general effective rational representative $A'$ of this ample
class.  Using a sufficiently high very ample multiple dilutes its
coefficients; Bertini on the same resolution ensures that adding it
preserves klt and gives it no component in common with $\Delta_0$.
Set
\[
 B'=\lambda G'+A',\qquad \Delta=\Delta_0+B'.
\]
Then $\Delta^+\leq\Delta_0^++\lambda G'+A'$ is klt.  Its positive
part is big because it contains the ample rational divisor $A'$.
Finally
\[
 K_Z+\Delta^+\sim_{\mathbb Q}\epsilon d^*W+\Delta^-
\]
is pseudoeffective.  BCHM big-boundary nonvanishing
\cite[Theorem~D]{BCHM2010} gives real-linear effectivity, and
Lemma~\ref{lem:rational-recovery} gives rational linear effectivity.
Pushing forward and dividing by $\epsilon$ proves the theorem with
\[
 C_A=\epsilon^{-1}d_*\Delta^-.
\]
This divisor is effective and supported on $A$: the only negative
terms introduced in \eqref{jet:delta-zero} were $K_{Z/Y}$, which is
exceptional, and $d^*A$.
\end{proof}

\subsection{The Iitaka base revisited}

We apply the jet construction to the Iitaka base already used in
Section~\ref{conv:iitaka}.  The algebraic return remains
Lemma~\ref{conv:return}; the new point is a fiber-intersection proof
that the original smooth metric is strictly positive in base directions.
We also record the stronger generic-fiber rank statement behind this
version of the argument.

Retain $X$ and $L=-K_X$ with its smooth semipositive metric from
Theorem~\ref{conv:fixed-error}.  Write the given effective integral
divisors as $N_m\sim mL-D$ for an unbounded set of positive integers
$m$, with $D$ pseudoeffective Cartier, and fix two exponents
$j<q$.  Use the Iitaka fibration of $N_q$ constructed in
Section~\ref{conv:iitaka}.  That construction settles the cases of
zero-dimensional and generically finite base.  In the remaining case
it gives
\[
 X\xleftarrow{\ \pi\ }S\xrightarrow{\ f\ }Y,
 \qquad 0<b=\dim Y<n=\dim X,
 \qquad u\pi^*N_q\sim f^*A_0+I,
\]
where $A_0$ is ample Cartier, $u$ is a positive integer, and $I\ge0$.  The varieties are
smooth projective, $f$ has connected fibers, and every source prime
mapping into base codimension at least two is $\pi$-exceptional.
Every ratio from a multiple of $N_q$ belongs to $\mathbb C(Y)$, which
is relatively algebraically closed in $\mathbb C(X)$.
These are precisely the model and field properties established there;
no flatness of the final smooth resolution is needed.

Put $E=K_{S/X}$ and $D_S=\pi^*N_q+E$.  Adding positive multiples of
$E$ does not change the sections of multiples of $\pi^*N_q$, because
$E$ is exceptional and $X$ is normal.  Thus $\kappa(D_S)=b$ and
$uD_S-f^*A_0$ has an effective representative.  In fact
\begin{equation}\label{jet:all-ranks}
 \operatorname{rank}f_*\mathcal O_S(\ell D_S)=1
 \qquad(\ell\in\mathbb Z_{>0}).
\end{equation}
For if this rank exceeded one, an ample base twist would yield two
sections of $\ell D_S+af^*A_0$ independent over $\mathbb C(Y)$.
Multiplication by the section of $a(uD_S-f^*A_0)$ places them in
$(\ell+au)D_S$.  Their ratio is transcendental over $\mathbb C(Y)$.
Other sections of a sufficiently divisible multiple recover $b$
algebraically independent base functions through $A_0$.  Products put
all these ratios into one multiple of $D_S$, contradicting
$\kappa(D_S)=b$.  The canonical sections show that the rank cannot be
zero.  At the generic point $\eta$ of $Y$, this gives
$h^0(S_\eta,E|_{S_\eta})=1$, and
$\pi^*N_q|_{S_\eta}$ has Iitaka dimension zero.

The interpolation and vertical-minimum construction of
Section~\ref{conv:iitaka}, with these indices, gives
\[
 R=\frac{\pi^*N_q-\pi^*N_j}{q-j}\sim_{\mathbb Q}\pi^*L,
 \qquad
 W=\sum_T\min_{V\mapsto T}
       \frac{\operatorname{ord}_V R}{\operatorname{ord}_V f^*T}\,T.
\]
Here $R$ has nonnegative horizontal part and $W$ is pseudoeffective.
The rank statement above also gives the interpolation directly:
for each given $m>q$, restrict the two effective divisors
$(m-j)N_q\sim(m-q)N_j+(q-j)N_m$ to the generic fiber.  Its
one-dimensional space of sections makes their horizontal divisors
equal.  The horizontal slope is therefore nonnegative as $m$ tends
to infinity.  The principal difference has vertical divisor, so its
restriction to the projective normal generic fiber is constant; hence
it is pulled back from $Y$.  Taking vertical minima in this equality
recovers the pseudoeffectivity of $W$ by the fixed-error limit used
above.

Call a base prime missed when every source prime above its generic
point is $\pi$-exceptional.  The model property also makes all other
components of its pullback exceptional.  By Lemma~\ref{conv:return},
it is enough to construct a weighted metric satisfying
Theorem~\ref{jet:base-lemma} on $-K_Y+A$, with $A$ supported on missed
primes.  We now verify the only new ingredient, horizontal strictness,
using intersections on the fibers.

\subsection{Horizontal positivity from fiber intersections}

Let $\widetilde L=\pi^*L$ and let $\alpha$ be its smooth semipositive
curvature form.  Since $D$ is pseudoeffective and
$u\pi^*N_q-f^*A_0$ has an effective representative,
$k\widetilde L-f^*A_0$ is pseudoeffective for some $k>0$.
Let $J$ be ample on $S$, put $d_0=n-b$, and choose the largest
$t\in\{0,\ldots,d_0\}$ such that
\[
 (\widetilde L^tJ^{d_0-t}\cdot S_y)>0
\]
on a smooth general fiber.  The choice exists because $t=0$ works.
These intersections are constant on the smooth base locus.
Semipositivity then makes the restriction of $\alpha$ have rank at
most $t$ at every point of every such fiber.

The pseudoeffective comparison and nefness give
\[
 ((k\widetilde L)^b\cdot\widetilde L^tJ^{d_0-t})
 \ge ((f^*A_0)^b\cdot\widetilde L^tJ^{d_0-t})>0.
\]
Indeed the difference of the $b$-th powers factors as the
pseudoeffective class $k\widetilde L-f^*A_0$ times a sum of products
of nef classes.  The final strict inequality is the projection
formula and the definition of $t$.  The smooth representative shows
that $\alpha$ has rank at least $b+t$ on a nonempty open set.

A vertical null vector for a nonnegative Hermitian form is null for
the full form.  The vertical null space has dimension at least
$d_0-t$, so the total rank is at most $b+t$ on the smooth locus.
On the preceding open set equality holds, and the full kernel is
vertical.  A smaller constant-rank patch, disjoint from $E$, therefore
satisfies
\[
 \alpha\ge c f^*\omega_Y
\]
for some $c>0$ and a positive base form $\omega_Y$.

The adjoint rank is one by \eqref{jet:all-ranks}.
Proposition~\ref{coa:volume-metric} now gives an integral effective
$A$ supported on missed primes and a semipositive metric on
$-K_Y+A$, smooth and strictly positive on an ordinary open set,
with weighted density in $L^{1+\gamma}_{\rm loc}$ for some
$\gamma>0$.  Theorem~\ref{jet:base-lemma} makes $W+C_A$ rationally
effective with $C_A\ge0$ supported on $A$.
Lemma~\ref{conv:return} applied to $R$ gives a section of a positive
multiple of $L$.  This completes the jet-based proof of
Theorem~\ref{conv:fixed-error}.

\section{A Fano-type model of the tensor fibration}
\label{fano:section}

The preceding effectivity arguments correct a divisor on a smooth base.
Here we instead contract the base until no correction is needed there.
The additional geometric conclusion is a Fano-type model of this
auxiliary base.  Its contracted divisors pull back to exceptional
divisors over the original variety, so a principal correction on the
model still produces a section on the original variety.

We give a Fano-type proof of the tensor conversion in
Corollary~\ref{conv:tensor-conversion}.  Its input is a smooth
connected projective $X$ with smoothly semipositive $L=-K_X$, and
nonzero sections of $(\Omega_X^1)^{\otimes p}\otimes mL$ for one
fixed $p\geq0$ and unbounded positive $m$.  The distinctive output
is the supported birational model of Theorem~\ref{fano:model}.
We first construct the auxiliary base and verify the hypotheses of
the toroidal volume theorem.  We then identify exactly which base
divisors its log terminal model can contract.  The final passage to
an anticanonical section uses the principal divisors themselves,
not only their numerical classes.

\subsection{A fibration tangent to the null directions}

The cases $\dim X=0$ and $p=0$ are immediate.  Suppose otherwise.
Write $\mathcal E=(\Omega_X^1)^{\otimes p}$ and apply the
determinant extraction in Lemma~\ref{conv:determinant}.
Retain its determinant line $M$, rather than its negative: there is a nonzero map
\[
 M\longrightarrow\bigwedge^s\mathcal E
\]
for some $s>0$, and effective integral divisors
\begin{equation}\label{fano:determinants}
 N_m\sim M+mL
\end{equation}
for an unbounded set of positive integers $m$.  The generic wedges
used in that lemma are regular because they take values in the
saturated subsheaf and its reflexive determinant.  This is the
determinant method of \cite[Lemma~4.1]{LP2018} and
\cite[Lemma~5.1]{LMPTX2023}, with the unbounded exponents supplied
here as part of the hypothesis.

Write $n=\dim X$, let $\alpha$ be the normalized curvature of $L$,
and let $\nu$ be its maximum rank.  If $\nu=n$, then $L$ is nef
and $L^n=\int_X\alpha^n>0$, so $L$ is big and conversion follows.
Fix an ample class $A$ and assume $\nu<n$.  We will use
\begin{equation}\label{fano:null-degree}
 N_m\cdot[\alpha]^\nu A^{n-1-\nu}=0.
\end{equation}
Here is a differential-geometric proof of this equality.  For each
$\epsilon>0$, Yau's prescribed-Ricci theorem~\cite{Yau1978} gives a
K\"ahler metric in $[\alpha]+\epsilon A$ whose Ricci form equals
the prescribed semipositive form $\alpha$, up to the fixed positive
normalization.  K\"ahler symmetry
identifies the mean curvature of the tangent bundle with its
nonnegative Ricci endomorphism.  Thus the mean curvature of every
cotangent tensor power, and of its exterior powers, is nonpositive.
Consequently
\[
 c_1(M)([\alpha]+\epsilon A)^{n-1}\leq0.
\]
Indeed, a positive degree would allow a metric on $M$ with constant
positive mean curvature, by the scalar Poisson equation.  The induced
mean curvature on
$\bigwedge^s\mathcal E\otimes M^{-1}$ would then be
negative definite, and the integrated Bochner formula would forbid
its nonzero section.  Since $\alpha^{\nu+1}=0$, the coefficient of
the lowest possible power $\epsilon^{n-1-\nu}$ is a positive
constant times $c_1(M)[\alpha]^\nu A^{n-1-\nu}$ and is nonpositive.
The same intersection for $N_m$ equals it by
\eqref{fano:determinants}, and is nonnegative by effectivity.  This
proves \eqref{fano:null-degree}.

Choose a finite sum $D$ of the $N_m$, allowing repetitions, whose
Iitaka dimension $d$ is maximal among such sums.  These dimensions
are integers in $\{0,\ldots,n\}$, so a maximum exists.  Choose
$q>0$ for which $|qD|$ has image dimension $d$.  On the open set
where $\alpha$ has rank $\nu$, the differential of this rational
map annihilates $\ker\alpha$.  To see this, each member of $|qD|$
has zero intersection with $[\alpha]^\nu A^{n-1-\nu}$.  A smooth
local component whose tangent hyperplane failed to contain
$\ker\alpha$ would retain rank $\nu$ for the restricted form,
and hence contribute positive mass to that intersection.  Away from
the base locus, a projective hyperplane detecting a nonzero
differential on $\ker\alpha$ would produce exactly such a smooth
local component.  Hyperplanes detect the differential of the map,
so the assertion follows.  In particular $d\leq\nu<n$.

Resolve the base ideal by $p_1:X_1\to X$, with $X_1$ smooth and the
resolution an isomorphism off that ideal, and take the Stein
factorization $f_1:X_1\to B$.  The normal projective base $B$ has
relatively algebraically closed function field in $\mathbb C(X)$.
There are an ample divisor $A_B$ and an effective fixed divisor $I$
with
\begin{equation}\label{fano:fixed-free}
 p_1^*(qD)\sim f_1^*A_B+I.
\end{equation}
Every $p_1$-exceptional prime lies in $\operatorname{Supp}I$, since
its center lies in the base ideal and that ideal has positive order
on the prime.

Fix $m_1<m_2$ among the determinant exponents and put
\[
 G=\frac{N_{m_2}-N_{m_1}}{m_2-m_1}\sim_{\mathbb Q}L,
 \qquad J=N_{m_1}-m_1G.
\]
On any compatible higher model $\mu:W\to X$, $f:W\to Y$, with
$Y$ birational over $B$, there are $\mathbb Q$-principal divisors
$T_m$ on $Y$ such that
\begin{equation}\label{fano:interpolation}
 \mu^*N_m=m\mu^*G+\mu^*J+f^*T_m\qquad(m>m_2).
\end{equation}
For the principal difference of
\[
 (m_2-m_1)N_m+(m-m_2)N_{m_1}
 \quad\hbox{and}\quad(m-m_1)N_{m_2}
\]
is a rational function on $B$.  Otherwise it is transcendental over
$\mathbb C(B)$, by relative algebraic closedness, and adjoining it
to the map for $|qD|$ increases the image dimension.  The joint map
is given by products of sections in the system for the permitted
finite sum $qD+(m-m_1)N_{m_2}$, contradicting maximality.
Dividing its principal divisor by $m_2-m_1$ proves
\eqref{fano:interpolation}.  If $d=0$, these functions are constant:
$N_m=mG+J$, so every coefficient of $G$ is nonnegative by letting
$m\to\infty$.  This already proves conversion.  We henceforth
have $0<d<n$.

We next establish the adjoint rank required for the volume metric.
For each integer $\ell>0$, the general fiber of $f_1$ satisfies
\begin{equation}\label{fano:fixed-rank}
 h^0(X_{1,b},\mathcal O_{X_1}(\ell I)|_{X_{1,b}})=1.
\end{equation}
If the dimension were larger, generic base change and a sufficiently
ample base twist would give sections of $\ell I+t f_1^*A_B$
whose ratio varies on the general fiber.  For $t\gg\ell$, multiply
them by the section of $(t-\ell)I$.  They become sections of
$t p_1^*(qD)$.  The same complete system recovers the base using the
sections from $|tA_B|$ multiplied by $s_I^t$, so its image dimension
exceeds $d$, a contradiction.  The Jacobian divisor $K_{X_1/X}$
is effective and bounded by some $\ell I$.  Its nonzero canonical
section and \eqref{fano:fixed-rank} therefore give rank one for its
full direct image.  Further smooth modifications do not change this
rank: their extra relative canonical divisors are effective and
exceptional, and their section pushforwards equal $\mathcal O$ by
normality.  Thus on any smooth resolution $\pi:Z\to X$,
$g:Z\to Y$ dominating these models,
\begin{equation}\label{fano:adjoint-rank}
 \operatorname{rank}g_*\mathcal O_Z(K_{Z/X})=1.
\end{equation}
The Jacobian section identifies
$g_*(K_{Z/Y}+\pi^*L)$ with $-K_Y$ on the smooth base-change locus.
This is an ordinary rank assertion, involving the entire adjoint
space.

\subsection{The integrable correction and the model it permits}

Apply birational weak toroidalization~\cite[Theorem~1.1]{ADK2013},
marking the nonisomorphism locus over $X$, and equidimensional
toroidal subdivision~\cite[Proposition~4.4]{AK2000}.  We obtain
$\mu:W\to X$, $f:W\to Y$ with $W$ normal projective, $Y$ smooth
projective, and $f$ surjective, equidimensional and toroidal for
strict toroidal embeddings.  All modifications are birational; no
reduced-fiber alteration is required.  The source can be singular,
and its resolution is used only for integration.  The effective
exceptional Weil divisor $K_{W/X}$ is supported in the toroidal
boundary.

Choose a global K\"ahler form $\omega_Y$ on $Y$.  By generic
smoothness and cohomology-and-base-change, there is a nonempty
Zariski open $U\subset Y$ on which the resolved map $g$ is smooth
and its adjoint direct image is the line identified above.  The
kernel calculation persists off the exceptional loci, since the map
from the base to the original image is generically finite.  On a
small patch above $U$, in the maximum-rank region and away from
the Jacobian divisor, it gives
\[
 \ker\pi^*\alpha\subset\ker dg,
 \qquad \pi^*\alpha\geq c\,g^*\omega_Y\quad(c>0).
\]
The second inequality follows from the first by semipositive linear algebra
and a lower bound for the nonzero eigenvalues on a smaller patch.
Together with \eqref{fano:adjoint-rank}, these verify all hypotheses
of the toroidal volume theorem in
\cite[Theorem~2.1]{CompanionG}.

Explicitly, for a base prime $P$ and primes $S\subset W$ dominating
it, put $a_S=\operatorname{ord}_S K_{W/X}$ and
$e_S=\operatorname{ord}_S f^*P$.  The theorem supplies
\begin{equation}\label{fano:correction}
 C=\sum_P\max\left\{0,\min_{S\mapsto P}
           \left(\frac{1+a_S}{e_S}-1\right)\right\}P
\end{equation}
and a semipositive metric on $\Lambda=-K_Y+C$ whose local squared
frame norms are integrable and whose curvature is smooth and
strictly positive on a nonempty ordinary open set.  Moreover $C$
is effective rational with finite support and $f^*C$ is entirely
exceptional over $X$.  The companion theorem proves the all-valuation
integrability and the extensions across divisors and codimension
two; these conclusions use the precise ordinary rank and patch
verified here.

In particular $\kappa(Y,C)=0$.  Indeed, a rational function with
poles bounded by a positive divisible multiple of $C$ pulls back
to a function with only $\mu$-exceptional poles.  It extends across
$X$ by normality and is constant because $X$ is connected projective.
Also there is an effective big rational divisor
\begin{equation}\label{fano:klt-representative}
 \Delta\sim_{\mathbb Q}\Lambda,
 \qquad (Y,\Delta)\ \hbox{klt}.
\end{equation}
Apply Corollary~\ref{cor:big-klt-representative} to the unweighted
integrable metric on $\Lambda$.  Its proof first subtracts a small
ample rational class while keeping a klt representative, then adds
a general representative of that ample class.  Thus the boundary
here is effective, big and klt, exactly as required by the birational
theorem below.

\begin{theorem}[A model with supported contractions]\label{fano:model}
Let $Y$ be smooth projective.  Suppose $C\geq0$ is a rational divisor
with $\kappa(Y,C)=0$, and $\Delta\geq0$ is big and rational with
$(Y,\Delta)$ klt and $K_Y+\Delta\sim_{\mathbb Q}C$.  Then there is
a birational contraction $\chi:Y\dashrightarrow Y_0$ to a projective
$\mathbb Q$-factorial variety of Fano type.  Every prime contracted
by $\chi$ lies in $\operatorname{Supp}C$.  Every pseudoeffective
rational divisor on $Y_0$ is $\mathbb Q$-linearly equivalent to an
effective rational divisor.
\end{theorem}

\begin{proof}
The big-boundary log terminal model theorem
\cite[Theorem~1.1]{BCHM2010} applies: the pair is klt, its boundary
is effective and big, and its adjoint is pseudoeffective since it
is rationally equivalent to $C\geq0$.  It gives a projective
$\mathbb Q$-factorial log terminal model $\chi:Y\dashrightarrow Y_0$.
The map extracts no divisors.  With $\Delta_0=\chi_*\Delta$, the
transformed pair is klt and $K_{Y_0}+\Delta_0$ is nef.  On a common
resolution $a:U\to Y$, $b:U\to Y_0$, its negativity property
\cite[Definitions~3.6.1 and~3.6.6]{BCHM2010} gives
\[
 a^*(K_Y+\Delta)=b^*(K_{Y_0}+\Delta_0)+E,
\]
where $E\geq0$ is $b$-exceptional and contains the strict transform
of every $\chi$-contracted prime.  Choose a $\mathbb Q$-principal
divisor $T$ with $K_Y+\Delta=C+T$.  Using the same rational
functions and coefficients downstairs, and then pulling back,
cancels $T$ on both sides.  Therefore, with $C_0=\chi_*C$,
\begin{equation}\label{fano:negative-equation}
 a^*C=b^*C_0+E.
\end{equation}
Since $Y_0$ is $\mathbb Q$-factorial, $C_0$ is rationally Cartier;
it is effective, so its pullback is effective.  The strict transform
of each contracted prime has positive coefficient on the right
of \eqref{fano:negative-equation}, and hence on the left.  That
prime must belong to $\operatorname{Supp}C$.

The boundary $\Delta_0$ remains big under the birational contraction.
The nef adjoint is semiample by big-boundary base-point-freeness
\cite[Corollary~3.9.2]{BCHM2010}.  The effective exceptional
term $E$ preserves the spaces of sections of all sufficiently
divisible adjoint multiples, by normality.  Thus
\[
 \kappa(Y_0,K_{Y_0}+\Delta_0)=\kappa(Y,C)=0.
\]
A semiample rational divisor with Iitaka dimension zero is rationally
linearly trivial: a base-point-free multiple defines a map to a
point and is therefore trivial.  Hence
$K_{Y_0}+\Delta_0\sim_{\mathbb Q}0$.

Write $\Delta_0\sim_{\mathbb Q}A_0+N_0$ with $A_0$ ample rational
and $N_0\geq0$.  For a sufficiently small rational $\lambda>0$, the
boundary $\Theta=(1-\lambda)\Delta_0+\lambda N_0$ is effective and
klt; this follows by checking the finitely many coefficients on a
common log resolution.  We obtain
\[
 -(K_{Y_0}+\Theta)\sim_{\mathbb Q}\lambda A_0,
\]
which is ample.  This proves that $Y_0$ is of Fano type.

Finally, let $B_0$ be a pseudoeffective rational divisor on $Y_0$.
For a small rational $t>0$, the class
$tB_0-(K_{Y_0}+\Theta)$ is ample.  Choose a general effective
rational representative $U_0$ of that class with sufficiently small
coefficients so that $(Y_0,\Theta+U_0)$ remains klt.  Its boundary
is effective and big, and its adjoint is rationally equivalent to
$tB_0$, which is pseudoeffective.  Big-boundary nonvanishing
\cite[Theorem~D]{BCHM2010} makes this adjoint real-linearly effective.
Lemma~\ref{lem:rational-recovery} makes it rationally linearly
effective.  Division by $t$ proves the assertion for $B_0$.
\end{proof}

\subsection{Returning the correction to the original variety}

Apply Theorem~\ref{fano:model} to \eqref{fano:klt-representative}.
We now finish this proof of Corollary~\ref{conv:tensor-conversion}.  For each base
prime $P$ set
\[
 g_P=\min_{S\mapsto P}\frac{\operatorname{coeff}_S\mu^*G}{e_S},
 \qquad
 j_P=\max_{S\mapsto P}\frac{\operatorname{coeff}_S\mu^*J}{e_S},
 \qquad B_g=\sum_Pg_PP,\quad B_j=\sum_Pj_PP.
\]
These are finite rational sums.  The horizontal coefficients of
$\mu^*G$ are nonnegative by \eqref{fano:interpolation} and
$m\to\infty$.  At a component attaining the minimum $g_P$, that
same equation and effectivity of $\mu^*N_m$ give
\[
 m g_P+j_P+\operatorname{coeff}_P T_m\geq0.
\]
Thus $mB_g+B_j+T_m\geq0$.  Push this inequality by strict transform
to $Y_0$.  Principal terms remain principal because $\chi$ extracts
no divisors.  Dividing numerical classes by $m$ proves that
$(B_g)_0$ is pseudoeffective.  Theorem~\ref{fano:model} supplies
a $\mathbb Q$-principal adjustment making $(B_g)_0$ effective.
Use its same rational functions and coefficients on $Y$, and call
the resulting principal divisor $T_*$.  Then $B_g+T_*$ is effective
outside the primes contracted by $\chi$.

The divisor $\mu^*G+f^*T_*$ has nonnegative coefficient at every
prime of $W$ that is not exceptional over $X$.  Horizontals were
already treated.  Above an uncontracted prime this follows from
the defining minimum $g_P$.  Above a contracted prime, every
component is exceptional over $X$, since that prime lies in
$\operatorname{Supp}C$ and $f^*C$ is exceptional.  Equidimensionality
leaves no other vertical primes.  Pushing forward to $X$ therefore
gives an effective rational divisor.  Its class is
$\mathbb Q$-linearly equivalent to $G\sim_{\mathbb Q}L$, because
the base principal adjustment pulls back and pushes forward as a
principal divisor.  Clearing denominators completes the proof.

The Fano-type conclusion belongs to $Y_0$, the model of an auxiliary
linear-system base.  The argument on $X$ uses exactly the supported
exceptional correction; it does not assert that $X$ itself is of
Fano type.

\section{Invariant sections from a curvature-null base}
\label{inv:section}

The ordinary conversion theorem does not control characters.  We now
construct the base from invariant ratios and retain an invariant
rational section throughout the return step.  The cohomological input
comes from the natural anticanonical index.  Here the natural
linearization on $-K_X=\det T_X$ is the action induced by the
differential of the action on $X$.  The subsequent determinant,
base, and metric constructions are local to the present argument.

\begin{proposition}\label{inv:forms}
Let $X$ be a compact K\"ahler manifold, let a compact real torus $T_c$
act holomorphically, and give $L=-K_X$ its natural linearization.
Suppose $L$ has a smooth semipositive metric and
\[
 I_{T_c}(0):=\sum_q(-1)^q\dim H^q(X,\mathcal O_X)^{T_c}\ne0.
\]
There are a fixed $p\ge0$ and unbounded positive integers $m$ with
$H^0(X,\Omega_X^p\otimes(m+1)L)^{T_c}\ne0$.
\end{proposition}
\begin{proof}
The invariant anticanonical index theorem \cite[Theorem~1.1]{CompanionI}
applies to the compact
complex manifold with its natural anticanonical linearization.  It
gives an integer $a>0$ and a polynomial $P$ such that
$I_{T_c}(m)=P(m)$ for every positive multiple of $a$, with the actual
value $P(0)=I_{T_c}(0)$.  In particular $P$ is not the zero polynomial,
so one degree $q$ supplies nonzero invariant cohomology at unbounded
positive divisible exponents.

Average a K\"ahler form and the logarithmic weights of the given
metric over $T_c$.  Hard Lefschetz with semipositive coefficients
\cite[Theorem~0.1]{DPS2001} gives the equivariant surjection
\[
 H^0(X,\Omega_X^{n-q}\otimes(m+1)L)
 \longrightarrow H^q(X,K_X+(m+1)L)=H^q(X,mL).
\]
The coefficient is smooth, so its multiplier ideal is trivial.
Compact averaging makes the map surjective on invariant vectors.
Set $p=n-q$.  The shift by one is necessary because the canonical
factor cancels one copy of the naturally linearized $L$.
\end{proof}

\begin{theorem}\label{inv:nonvanishing}
Let $X$ be smooth connected projective with
$\chi(X,\mathcal O_X)\ne0$.  Let an algebraic torus $T$ act on $X$,
and suppose $L=-K_X$ is smoothly semipositive.  For the natural
anticanonical linearization,
$H^0(X,mL)^T\ne0$ for some positive integer $m$.
\end{theorem}

A connected torus acts trivially on ordinary cohomology by isotopy and
hence on $H^q(X,\mathcal O_X)$ by Hodge theory.  Thus
$\chi(X,\mathcal O_X)$ is exactly the zero value in
Proposition~\ref{inv:forms}.  The point case is immediate.  For
$n=\dim X>0$, average the metric over the maximal compact torus and
write its curvature as $\theta\ge0$, normalized to represent $c_1(L)$.
Let $\nu$ be its maximum rank and let $\omega_0$ be a K\"ahler form.
The invariant forms supplied by Proposition~\ref{inv:forms} give the
answer immediately if their degree is zero.  Otherwise determinant
extraction as in Lemma~\ref{conv:determinant} gives a linearized line
$A$, a nonzero map $A\to\bigwedge^e\Omega_X^p$, and invariant sections
\begin{equation}\label{inv:determinant-sequence}
 0\ne\sigma_j\in H^0(X,A+m_jL)^T,
 \qquad 0<m_1<m_2<\cdots.
\end{equation}
We will use an intersection consequence of the determinant map, rather
than replace this construction by the ordinary conversion theorem.

\begin{lemma}\label{inv:null-slope}
If $\nu<n$, the determinant line just constructed satisfies
\[
 c_1(A)\cdot[\theta]^\nu[\omega_0]^{n-1-\nu}\le0.
\]
\end{lemma}
\begin{proof}
For each $t>0$, Yau's prescribed-Ricci theorem \cite{Yau1978} gives
a K\"ahler metric in $[\omega_0+t\theta]$ with Ricci form
$2\pi\theta$.  The induced mean curvature on
$\bigwedge^e\Omega_X^p$ is negative semidefinite.  If $A$ had
positive degree in this K\"ahler class, solving the scalar Laplace
equation would give it a metric with constant positive curvature
trace.  The induced mean curvature on
$\operatorname{Hom}(A,\bigwedge^e\Omega_X^p)$ would then be negative
definite, contradicting the nonzero holomorphic map by the integrated
Bochner formula.  Hence
$c_1(A)\cdot[\omega_0+t\theta]^{n-1}\le0$ for every $t>0$.
Because $\theta^{\nu+1}=0$, the leading coefficient as $t\to\infty$
is the asserted nonpositive intersection, up to a positive binomial
factor.
\end{proof}

\subsection{The invariant ratio field}

Put $B_j=A+m_jL$.  Form the subfield generated by ratios of invariant
sections of the same line bundle $\sum_j k_jB_j$, where $k_j\ge0$
are integers of finite support.  Let $F$ be its relative algebraic
closure in $\mathbb C(X)$.  Both fields are finitely generated:
choose a transcendence basis for the smaller field; the algebraic
closure of the associated purely transcendental field in the finitely
generated extension $\mathbb C(X)$ is finite, and the intervening
algebraic fields are finite as well.  The field $F$ is fixed by $T$.
Indeed an element algebraic over the invariant subfield has finite
orbit among the roots of a polynomial, and connectedness makes this
orbit a point.

Finitely many ratios can be put over a common denominator by
multiplication of sections.  Thus one bundle $B=\sum k_jB_j$ and
invariant sections of $B$ define an image $Y_0$ with the smaller
function field.  Normalize in $F$, resolve the base, and obtain the
invariant rational map $X\dashrightarrow Y$, with
$\mathbb C(Y)=F$ and geometrically integral generic fiber.
Let $H_0$ be the pullback of the image polarization; it is big and
base point free.  On a smooth equivariant graph resolution
\[
 X\xleftarrow{\ \pi\ }Z\xrightarrow{\ g\ }Y
\]
there is a nonzero invariant section of
$\pi^*B-g^*H_0$: divide the pulled-back defining sections by the
homogeneous coordinates on the image.  Hence for a chosen ample $H$
on $Y$ there is an integer $b>0$ and an invariant section of
$b\pi^*B-g^*H$.

If $F=\mathbb C$, consider
\begin{equation}\label{inv:ratio-identity}
 r_j=\frac{\sigma_j^{m_2-m_1}\sigma_1^{m_j-m_2}}
                 {\sigma_2^{m_j-m_1}},\qquad j>2.
\end{equation}
Its numerator and denominator are invariant sections of the same
multiple of $B_2$, so $r_j$ is constant.  Divisor comparison and
$m_j\to\infty$ show that
$\operatorname{div}(\sigma_2/\sigma_1)\ge0$.  This is the desired
invariant section of $(m_2-m_1)L$.  We may therefore suppose
$\dim Y>0$.

There is a horizontal strictness patch on $Z$.  If $\nu=n$, use any
maximum-rank point away from the exceptional and critical loci.
If $\nu<n$, Lemma~\ref{inv:null-slope} gives
$c_1(bB)\cdot[\theta]^\nu[\omega_0]^{n-1-\nu}\le0$.
Since $b\pi^*B-g^*H$ is effective, and all the forms involved are
semipositive, it follows that
\[
 \int_Z g^*\omega_H\wedge
           \pi^*(\theta^\nu\wedge\omega_0^{n-1-\nu})=0.
\]
At a maximum-rank point where $\pi$ is an isomorphism, nonnegative
linear algebra implies $\ker\pi^*\theta\subset\ker dg$.
Choose such a point over the smooth base-change locus.  Shrinking a
constant-rank neighborhood gives
\begin{equation}\label{inv:horizontal-patch}
 \pi^*\theta\ge c g^*\omega_H,\qquad c>0,
\end{equation}
on a nonempty ordinary patch, disjoint from the Jacobian divisor.
Rank zero is impossible here when $\dim Y>0$.

On this equivariant model, with the trivial action on $Y$, the invariant
part of $g_*\mathcal O_Z(K_{Z/X})$ has rank one.  The natural Jacobian
section gives rank at least one.  If it were larger, global generation
after a sufficiently ample base twist would give two invariant
sections independent on the generic fiber.  Multiplying them by a
power of the invariant section of $b\pi^*B-g^*H$ puts them in
$K_{Z/X}+k b\pi^*B$.  Birational pushforward, using
$\pi_*\mathcal O_Z(K_{Z/X})=\mathcal O_X$, gives a ratio of
invariant sections of $kbB$ not in $F$, a contradiction.
This is invariant rank one; no assertion about the rank of the full
direct image has been made.

We have obtained the analytic data: smooth invariant coefficients,
invariant adjoint rank one, the natural nonvanishing Jacobian generator,
and the horizontal patch \eqref{inv:horizontal-patch}.  To obtain
unweighted integrability with exceptional support, we now use a
separate normal toroidal model.  Equivariance of that model is not
required because the meromorphic data on $X$ are already invariant.

\subsection{A klt model and its unweighted volume}
\label{inv:toroidal-volume}

By weak toroidalization and equidimensional subdivision
\cite{ADK2013,AK2000}, prepare
\[
 p:W\to X,\qquad f:W\to Y
\]
with $W$ normal and simplicial toroidal, $Y$ smooth, $f$
equidimensional and toroidal, and the $p$-exceptional locus in the
boundary.  The base may be modified birationally, and the smooth
equivariant graph resolution above is then replaced accordingly.
For clarity, the simplicial condition means that the primitive rays of each local
toric cone are linearly independent.  Such toric charts are finite
abelian quotients of smooth toric charts.  The simplicial refinement
can be made projectively using
generic heights on the existing rays.  Since the base cones are
simplicial and source rays map to base rays or zero, each subdivided
cone still maps onto a face of a base cone.  Equidimensionality is
preserved.  Simplicial toroidal charts have quotient singularities;
in particular $W$ is $\mathbb Q$-factorial and klt.

Put $D_j=\operatorname{div}(p^*\sigma_j)$ and
$V=(D_2-D_1)/(m_2-m_1)\sim_{\mathbb Q}p^*L$.
Equation~\eqref{inv:ratio-identity} gives
\[
 D_j=D_1+(m_j-m_1)V+
            \frac{f^*\operatorname{div}_Y(r_j)}{m_2-m_1}.
\]
Thus $V$ has nonnegative horizontal part.  Every vertical prime of $W$
dominates a base prime.  Its minimum divisor
\[
 M=\sum_Q\left(\min_{D\mapsto Q}
             \frac{\operatorname{ord}_D V}{\operatorname{ord}_D f^*Q}
        \right)Q
\]
is pseudoeffective by the coefficientwise finite-error argument
following \eqref{conv:affine-divisors}.  On $W$ one has the actual inequality
$V\ge f^*M$.

The smooth metric of $L$ defines a positive volume $\mu$ on $X$.
On a smooth base-change open, let $\rho$ be the coordinate density of
$g_*\pi^*\mu$.  This measure can equally be computed on $W_{\rm reg}$;
birational changes affect only sets of measure zero.  Its density is
the squared norm of the Jacobian adjoint generator divided by the base
canonical frame.  Compact averaging on the smooth direct image is a
holomorphic orthogonal projection onto the invariant line.  Its norm
is therefore the quotient norm.  For positive-dimensional fibers,
Lemma~\ref{coa:curvature} gives
\[
 i\partial\bar\partial(-\log\rho)\ge0,
\]
with strictness on a nonempty ordinary open by
\eqref{inv:horizontal-patch}.  If the relative dimension is zero,
connected general fibers make $g$ birational.  Over its isomorphism
open the integral metric is simply the transported coefficient metric;
its curvature gives the same assertion directly, including strictness
on the patch.  This use of a smooth coefficient and an orthogonal
quotient does not assert positivity of an arbitrary subbundle of a
singular direct image.

Write $E_W=K_{W/X}\ge0$.  For a base prime $Q$, set
\[
 a_D=\operatorname{ord}_D E_W,\qquad
 e_D=\operatorname{ord}_D f^*Q,\qquad
 b_Q=\max\left\{0,\min_{D\mapsto Q}
                         \frac{1+a_D}{e_D}-1\right\},
 \quad F_0=\sum_Q b_QQ.
\]
This effective rational divisor has finite support, and
$f^*F_0\le E_W$.  In fact, if $b_Q>0$, then
$e_Db_Q\le1+a_D-e_D\le a_D$ for every $D\mapsto Q$.
If a prime above $Q$ is nonexceptional over $X$, its $a_D$ is zero
and $b_Q=0$.  Hence $f^*F_0$ is exceptional over $X$.

On $-K_Y+F_0$ consider the weight
\[
 \Phi=-\log\rho+\sum_Qb_Q\log|q_Q|^2,
\]
initially off the indicated divisors, with $q_Q$ a local equation of
$Q$.  At a general point of a prime $D$ attaining the minimum, choose
coordinates with $f=(z_1^{e_D},z_2,\ldots,z_{\dim Y})$.
The pulled-back volume density is comparable to
$|z_1|^{2a_D}$ times a smooth positive density.  Integration of this
chart gives
\[
 \rho(y)\ge c|q_Q(y)|^{2((1+a_D)/e_D-1)}.
\]
Consequently $\Phi$ is locally bounded above at a general point of
every omitted divisor.  Plurisubharmonic extension, first across these
divisors and then across codimension two, gives a semipositive metric
on the rational line $-K_Y+F_0$, still smooth and strictly positive
on an ordinary open.

Its squared frame norm is unweighted locally integrable.  To check all
valuations, take a log resolution $h:W'\to W$ of the relevant data.
Integrating $e^{-\Phi}$ on the base is the same as integrating the
pulled-back volume divided by the local divisor factors of $f^*F_0$.
The exponent divisor upstairs is
\[
 K_{W'}-(ph)^*K_X-h^*f^*F_0
  =K_{W'/W}+h^*(E_W-f^*F_0).
\]
The first summand has coefficients greater than $-1$ because $W$ is
klt, and the second is effective because $E_W-f^*F_0$ is effective
and rationally Cartier.  On a normal-crossings resolution all
coefficients are therefore greater than $-1$, which is exactly the
local Jacobian integrability test.  This includes exceptional divisors
above higher-codimension centers.

Apply Corollary~\ref{conv:unweighted-effectivity} to the pseudoeffective
$M$.  It gives $F_0+\epsilon M\sim_{\mathbb Q}G\ge0$ for some
rational $\epsilon>0$.  Pullback and the inequality $V\ge f^*M$
make $V$ effective after a principal correction from $Y$ and a divisor
exceptional over $X$.  Equivalently apply Lemma~\ref{conv:return} on
a common resolution.  The underlying rational section is a power of
$\sigma_2/\sigma_1$ times a base function, so it is invariant.
Its extension on $X$ proves Theorem~\ref{inv:nonvanishing}.

\section{A maximal invariant base and signed adjoint metrics}
\label{inv:dominated-alternative}

There is another invariant construction when
$H^q(X,\mathcal O_X)=0$ for every $q>0$.  The maximal base is chosen
from all invariant rational maps whose polarizations are
pseudoeffectively dominated by $L$.  Its distinguishing output is a
sequence of signed boundaries, one for every positive adjoint twist.
The limit of those boundaries proves the pseudoeffectivity needed for
nonvanishing; the metric for a single twist would not suffice.

More precisely, if $X$ is smooth connected projective over $\mathbb C$
with $H^q(X,\mathcal O_X)=0$ for all $q>0$, a compact real torus
$T_c$ acts holomorphically, and the naturally linearized $L=-K_X$
is smoothly semipositive, this construction gives
$H^0(X,mL)^{T_c}\ne0$ for some $m>0$.  Since
$H^1(X,\mathcal O_X)=0$, the torus preserves the
isomorphism class of a very ample polarization and acts in the
projective linear stabilizer of its image.  A compact torus there lies
in an algebraic torus after conjugation; taking its Zariski closure
therefore permits equivariant projective graph resolutions.  Averaging
metric weights gives a smooth invariant coefficient.  The zero-index
value is one.  Proposition~\ref{inv:forms} and
Lemma~\ref{conv:determinant} supply a linearized $A$ with $-A$
pseudoeffective and sections as in
\eqref{inv:determinant-sequence}, unless a scalar invariant section
has already been obtained.

\subsection{Maximality and all positive adjoint ranks}

Consider invariant dominant rational maps to projective bases with
trivial action, subject on a resolution $\pi:Z\to X$, $g:Z\to Y$
to
\begin{equation}\label{inv:domination}
 \pi^*L-\delta g^*H\quad\hbox{pseudoeffective}
\end{equation}
for some ample $H$ and rational $\delta>0$.  Include the point map.
This condition is unchanged by replacing the resolution.  It persists
under birational modification of the base or its finite Stein
refinement: the pullback of an ample divisor under a generically finite
map is big and dominates a small ample class.  It also persists for the
joint image of two such maps, by addition and a product polarization.
Choose a map of maximal base dimension, take the relative algebraic
closure of its field in $\mathbb C(X)$, and resolve the base.  The
resulting field is still invariant, as in the finite-orbit argument
above.  Every other admissible invariant ratio belongs to it: its
adjoining cannot increase dimension, and the field is already
relatively algebraically closed.

Fix two exponents $a<b$ in the determinant sequence, set $m=b-a$, and
put $t=\sigma_b/\sigma_a$, an invariant rational section of $mL$.
For every larger exponent $i$, the ratio
\[
 \frac{\sigma_b^{i-a}}{\sigma_i^{b-a}\sigma_a^{i-b}}
\]
belongs to $\mathbb C(Y)$.  It is a pencil in a bundle $jL+dA$
with $j,d>0$, and $-A$ pseudoeffective makes that pencil satisfy
\eqref{inv:domination}.  Thus the horizontal orders of the
$\sigma_i$ vary affinely in $i$.  Their nonnegativity at unbounded
$i$ makes $D_t=\operatorname{div}(\pi^*t)$ effective horizontally.
If $Y$ is a point, this gives the desired section directly.

On a smooth equivariant graph resolution set $E=K_{Z/X}$.
For every integer $k\ge1$ the invariant generic adjoint rank is
\begin{equation}\label{inv:all-ranks}
 \operatorname{rank}\bigl(g_*\mathcal O_Z(E+km\pi^*L)\bigr)^{T_c}=1.
\end{equation}
The rational section $s_Et^k$ is regular on the generic fiber, hence
supplies a nonzero invariant element.  If the rank exceeded one,
global generation after twisting by $jg^*H$ would give two invariant
sections in $E+km\pi^*L+jg^*H$ whose ratio is not a base function.
Their effective divisors push to a pencil in
$D_0=kmL+j\pi_*g^*H$ on $X$; exceptionality removes $E$.
Pushing \eqref{inv:domination} forward gives an integer $b>0$ for
which $bL-D_0$ is pseudoeffective.  Resolve the pushed-down pencil by
$\sigma:U\to X$ and $h:U\to\mathbb P^1$.  Its fixed part $F$ is
effective and
$\sigma^*D_0\sim h^*\mathcal O_{\mathbb P^1}(1)+F$.  Consequently
\[
 b\sigma^*L-h^*\mathcal O_{\mathbb P^1}(1)
 \sim_{\mathbb Q}\sigma^*(bL-D_0)+F
 \quad\hbox{is pseudoeffective}.
\]
This is exactly the domination premise for the new invariant pencil.
Its ratio is not a base function, contradicting maximality.

We use a normal $\mathbb Q$-factorial simplicial toroidal
equidimensional model $p:W\to X$, $f:W\to Y$, as above, for the
divisorial calculations.  It need not carry the action.  All positivity
on invariant summands will instead be proved on the smooth equivariant
model, with base change and the ranks \eqref{inv:all-ranks}.

\subsection{The signed boundaries and their integrals}

For a prime $D\mapsto Q$ on $W$, put
\[
 d_D=\operatorname{ord}_D\operatorname{div}(p^*t),\qquad
 a_D=1+\operatorname{ord}_D K_{W/X},\qquad
 e_D=\operatorname{ord}_D f^*Q.
\]
Here $a_D\ge1$, and $a_D=1$ for nonexceptional primes.  Define
\[
 \gamma_Q=\min_{D\mapsto Q}\frac{d_D}{e_D},\qquad
 \lambda_Q(k)=\min_{D\mapsto Q}\frac{kd_D+a_D}{e_D}.
\]
The first minimum records the vertical slope of $t$; the second
includes the Jacobian order relevant to integration.  We choose the
coefficient $\beta_Q(k)$ of the base correction to satisfy
\[
 \lambda_Q(k)-1\le\beta_Q(k)<\lambda_Q(k),\qquad
 k\gamma_Q\le\beta_Q(k),\qquad
 \beta_Q(k)\le\frac{kd_D}{e_D}
        \quad(D\text{ nonexceptional}).
\]
These inequalities have different roles.  Near general points of
$D$ and $Q$, the local map $y_1=z_1^{e_D}$ turns a volume density
of order $2(kd_D+a_D-1)$ into a contribution to the base density
with exponent $2((kd_D+a_D)/e_D-1)$.  A prime attaining the minimum
therefore gives a lower bound with exponent $2(\lambda_Q(k)-1)$.
Dividing the base density by $|y_1|^{2\beta_Q(k)}$ adds
$\beta_Q(k)\log|y_1|^2$ to its negative logarithm.  The first
lower bound makes this corrected weight locally bounded above at $Q$.
The divided volume upstairs has exponent
$2(kd_D+a_D-1-e_D\beta_Q(k))$; the strict upper bound makes it
integrable at each prime over $Q$.  The toroidal calculation below
will check the remaining valuations.  The final upper bound ensures
that the section returned to $X$ has no pole at a nonexceptional
prime.  The comparison with $k\gamma_Q$ will be used after taking
the limit of the correction divisors.

All four requirements are met by
\[
 \beta_Q(k)=\max\{k\gamma_Q,\lambda_Q(k)-1\},\qquad
 B_k=\sum_Q\beta_Q(k)Q.
\]
The supports lie in one finite set.  Since every $a_D/e_D>0$,
both terms in the maximum are strictly below $\lambda_Q(k)$.
For a nonexceptional $D$, both $k\gamma_Q$ and
$\lambda_Q(k)-1\le kd_D/e_D+1/e_D-1$ are at most $kd_D/e_D$.
Finally, $\lambda_Q(k)-k\gamma_Q$ is bounded independently of $k$,
because the minimum is over a fixed finite set.  Thus
\begin{equation}\label{inv:beta-bounds}
 \beta_Q(k)<\lambda_Q(k),\qquad
 \beta_Q(k)\le\frac{kd_D}{e_D}\quad(D\text{ nonexceptional}),
 \qquad \frac{B_k}{k}\longrightarrow
 B_\infty:=\sum_Q\gamma_QQ.
\end{equation}
The chosen lower bound by $k\gamma_Q$ also gives $B_1\ge B_\infty$.
After the metrics show that $-K_Y+B_k$ is pseudoeffective for every
$k$, the displayed limit will imply pseudoeffectivity of $B_\infty$
and hence of $B_1$.

The pair
\begin{equation}\label{inv:sub-klt}
 (W,f^*B_k-k\operatorname{div}(p^*t)-K_{W/X})
\end{equation}
is sub-klt.  Horizontal coefficients are nonpositive.  A vertical
coefficient equals $e_D\beta_Q(k)-kd_D-(a_D-1)<1$.
If $e_D=1$, both terms in the maximum defining $\beta_Q(k)$ give
this coefficient at most zero.  Thus the positive part is supported on
the toroidal boundary.  On a simplicial toroidal chart the discrepancy
function has positive values $1-b_D$ on the boundary rays, and every
new ray has a positive linear combination of those values.  Hence an
effective boundary with coefficients below one is klt there.
Subtracting an effective rational divisor only improves discrepancies.
This proves \eqref{inv:sub-klt}, including all divisorial valuations.

Let $d\mu$ be the smooth anticanonical volume on $X$.  On the smooth
base-change locus the invariant rank-one generator of
$K_{Z/Y}+(km+1)\pi^*L$ is the relative form associated to
$s_Et^k$.  The squared norm in its coordinate frame is
\[
 \rho_k=\frac{g_*\pi^*(\|t\|^{2k}d\mu)}{dV_Y}.
\]
Smooth direct-image positivity applies to the smooth invariant
coefficient $(km+1)\pi^*L$.  Compact averaging is a holomorphic
orthogonal projection, so the invariant line has the quotient metric.
The squared frame norm on $G_k=-K_Y+B_k$ is
\[
 H_k=\rho_k\prod_Q|q_Q|^{-2\beta_Q(k)}.
\]
Near a prime attaining the minimum in $\lambda_Q(k)$, the local map
has form $y_1=z_1^{e_D}$ and the numerator volume is comparable to
$|z_1|^{2(kd_D+a_D-1)}$ times smooth volume.  Therefore
\[
 \rho_k(y)\ge c|q_Q(y)|^{2(\lambda_Q(k)-1)}.
\]
Since $\beta_Q(k)\ge\lambda_Q(k)-1$, the weight $-\log H_k$
is locally bounded above at general points of all missing divisors.
It extends plurisubharmonically there and across codimension two.
Thus every $G_k$ is pseudoeffective.

Moreover $H_k$ is unweighted locally integrable.  On a log resolution
$h:W'\to W$, the divisor of poles in the integrand is
\[
 h^*(f^*B_k-k\operatorname{div}(p^*t))-K_{W'/X}
 =h^*(f^*B_k-k\operatorname{div}(p^*t)-K_{W/X})-K_{W'/W}.
\]
Its coefficients are below one by \eqref{inv:sub-klt}.  The
normal-crossings Jacobian test and properness prove integrability
on the base.  This signed boundary calculation is separate from the
effective exceptional correction in Section~\ref{inv:toroidal-volume}.

\subsection{A global strict metric and the final section}

We give the strict-metric construction with its actual sheaf and ideal
conditions.  From \eqref{inv:domination} choose a singular metric on
$\pi^*L$ with curvature at least $\delta g^*\omega_H$.
Average its global weight difference $u$ from the smooth invariant
metric over $T_c$, in the sense of distributions.  This preserves the
quasiplurisubharmonic lower bound.  On $(m+1)\pi^*L$, add
$\epsilon u$ to the smooth coefficient weight.  For small
$\epsilon>0$ the resulting metric $h_\epsilon$ satisfies
\[
 \Theta_{h_\epsilon}\ge\epsilon\delta g^*\omega_H,
 \qquad\mathcal J(h_\epsilon)=\mathcal O_Z.
\]
Small-exponent integrability on compact $Z$ proves the ideal assertion.
Consequently the singular adjoint direct-image theorem
\cite[Theorem~3.3.5]{PT2018} acts on the full ordinary direct-image
sheaf.  Let $U\subset Y$ be a smooth base-change open on which this
sheaf is locally free.  Fubini makes the fiber integral norm finite
almost everywhere on $U$.

On $U$, averaging over the compact torus defines a holomorphic
projection $P$ onto the invariant summand.  Invariance of
$h_\epsilon$ makes $P$ orthogonal almost everywhere, where the integral
norm is finite.  The dual of the quotient embeds holomorphically in the
dual full bundle, and its dual norm is the restricted norm.  Logarithms
of these dual norms are plurisubharmonic by direct-image positivity.
They therefore define the semipositive quotient metric; its rank is
one by \eqref{inv:all-ranks}.  Subtracting a base potential before
applying this argument retains the lower curvature bound
$\epsilon\delta\omega_H$.  This proves positivity for this specified
quotient, without a claim about arbitrary singular invariant subbundles.

For $k=1$ the integral norm is obtained by inserting $e^{-\epsilon u}$
in the density $\rho_1$.  Since $u$ is locally bounded above, the
lower density estimates remain valid up to constants.  The same
extension across divisors, and then codimension two, gives a metric on
$G_1=-K_Y+B_1$ with a global strict curvature lower bound.
Together with the already constructed integrable semipositive metric,
strong openness and H\"older's inequality give a globally strict
metric on $G_1$ whose squared frame norms remain locally integrable.
Indeed take a sufficiently small positive fraction of the strict
weight; the integrable metric has some $L^{1+\eta}$ room by strong
openness, and the new singular weight has small-exponent integrability.

It remains to verify the algebraic premise for the adjoint class.
Every $G_k/k$ is pseudoeffective, and
$G_k/k=-K_Y/k+B_k/k\to B_\infty$.
The closedness of the pseudoeffective cone gives $B_\infty$
pseudoeffective.  Since $B_1\ge B_\infty$,
$B_1$ is pseudoeffective as well.  Choose a small ample rational $A'$
whose smooth curvature can be subtracted from the global strict lower
bound on $G_1$.  The resulting metric on $G_1-A'$ remains globally
strict and unweighted integrable.  Apply
Lemma~\ref{lem:weighted-resolution} with $J=0$ and a sufficiently
large divisible $r>1$.  On its log resolution, the fixed part divided
by $r$, minus the relative canonical divisor, has coefficients below
one.  A general member of the free part, divided by $r$, meets this
support transversely and has coefficient $1/r<1$.  It therefore gives
an effective klt $\Delta_0\sim_{\mathbb Q}G_1-A'$.
Add a general effective rational representative of $A'$ with small
coefficients, transverse on a common log resolution.  The resulting
$\Delta\sim_{\mathbb Q}G_1$ is effective, big, and klt.
Only the global strict bound and unweighted integrability are used
here; the mixed metric need not be smooth on an open set.
Now $K_Y+\Delta\sim_{\mathbb Q}B_1$ is pseudoeffective, so
big-boundary nonvanishing \cite[Theorem~D]{BCHM2010} gives
\[
 \operatorname{div}(h)+NB_1\ge0
\]
for some $N>0$ and $h\in\mathbb C(Y)^*$, after applying Lemma~\ref{lem:rational-recovery}.
The rational section $t^Nf^*h$ has no pole at horizontal primes.  At a
vertical prime nonexceptional over $X$ this follows from
$\beta_Q(1)\le d_D/e_D$ in \eqref{inv:beta-bounds}.
Equidimensionality handles every such prime on $W$.
It therefore extends to a nonzero section of $NmL$ on $X$.
Both $t$ and the base function are invariant, proving invariant
nonvanishing under the stated cohomology-vanishing hypothesis by this
second construction.

\section{Compact isometric monodromy and descent}
\label{inv:isometric-descent}

Let $X$ be a smooth connected projective complex variety whose
anticanonical bundle has a smooth Hermitian metric of semipositive
Chern curvature. The finite-cover structure theorem supplies a
compact factor to which the preceding invariant results apply.
There are two legitimate orders of operation: descend
invariant scalar sections, or descend twisted forms before applying an
ordinary conversion theorem.  We give the transport in both cases,
since the second does not assume an equivariant conclusion from
Theorem~\ref{conv:fixed-error}.

The finite-cover structure theorem
\cite[Theorem~2.1]{CompanionH}
supplies a
connected finite \'etale cover $\widehat X\to X$ whose universal cover
has the holomorphic isometric splitting
\[
 \widetilde X=\mathbb C^a\times R\times F.
\]
Here $R$ is a product of compact simply connected Ricci-flat factors,
$F$ is a product of the remaining compact projective irreducible
factors, and the deck group projects to a full translation lattice on
$\mathbb C^a$.  Its action on the compact factors lies in a compact
connected torus of holomorphic isometries, and the flat and Ricci-flat
factors have invariant nowhere-vanishing canonical frames.  Moreover
$H^q(F,\mathcal O_F)=0$ for $q>0$, and $L_F=-K_F$ is smoothly
semipositive.  Empty products and point factors are allowed.
These geometric and canonical-frame properties are supplied by the
cited theorem and its proof. The separate anticanonical nonvanishing
theorem of \cite[Theorem~1.1]{CompanionH} is not used here.

\begin{remark}\label{inv:tensor-holonomy}
One can strengthen the form-vanishing part of the structure description
to positive covariant tensors by a short local holonomy argument.
Let $S$ be one compact simply connected nonflat irreducible K\"ahler
factor, with nonnegative Ricci curvature, and let
$H\subset U(T_xS)$ be its holonomy group.  It is connected, and its
complex tangent representation is irreducible by de Rham
irreducibility.  If its determinant character is trivial, the canonical
connection is flat and $S$ is Ricci-flat; simple connectivity then
gives a parallel canonical frame.

Otherwise the differential of the determinant character is nonzero.
In the compact Lie algebra
$\mathfrak h=\mathfrak z\oplus[\mathfrak h,\mathfrak h]$, trace
vanishes on commutators.  Some central element therefore has nonzero
trace.  Schur's lemma makes that element a nonzero imaginary scalar,
whose exponential supplies the full scalar circle in $H$.
On $(T_x^*S)^{\otimes p}$ this circle has weight $-p$, so it has no
invariant vector for every $p>0$.  The K\"ahler Bochner formula with
nonnegative Ricci curvature makes every holomorphic covariant tensor
parallel.  Hence
\[
 H^0(S,(\Omega_S^1)^{\otimes p})=0\qquad(p>0).
\]
This argument concerns all covariant tensor powers, not only exterior
forms.  It is independent of the descent map below, which already
works for the alternating forms supplied by hard Lefschetz.
\end{remark}

For scalar descent, apply Theorem~\ref{inv:nonvanishing} to $F$,
or the compact-torus construction of
Section~\ref{inv:dominated-alternative}.  The compact action is
contained in an algebraic torus: $H^1(F,\mathcal O_F)=0$ preserves
a very ample polarization, and a compact abelian projective linear
group is conjugate into a diagonal torus.  Thus a positive power of
$L_F$ has a section fixed by the projected deck action.
Pull it to $\widetilde X$ and multiply its coefficient by the matching
inverse powers of the invariant canonical frames on $\mathbb C^a$
and $R$.  The resulting anticanonical section is deck-invariant and
nonzero, so it descends to $\widehat X$.  If $F$ is a point, begin
with its constant section.  The finite \'etale section norm then gives
a nonzero anticanonical multiple on $X$.

For the second order of operation, Proposition~\ref{inv:forms} gives a
fixed $p$ and invariant nonzero sections
\[
 u_m\in H^0(F,\Omega_F^p\otimes L_F^{m+1})
\]
at unbounded positive $m$, since the actual zero index equals one.
The pullback map for differential forms on the product is injective.
Tensoring the coefficient with the $(m+1)$-st inverse powers of the
canonical frames on the other factors gives a nonzero form with
coefficient $K_{\widetilde X}^{-(m+1)}$.  It is deck-invariant:
translations preserve the flat frame, compact monodromy preserves the
Ricci-flat frame, and $u_m$ is invariant on $F$.  It therefore descends
to a nonzero section of
\[
 \Omega_{\widehat X}^p\otimes(-(m+1)K_{\widehat X}).
\]
These holomorphic sections are algebraic because $\widehat X$ is
projective.  If $p=0$, we have a scalar section already.
Otherwise Corollary~\ref{conv:tensor-conversion} converts the descended
forms on the actual projective cover to an ordinary anticanonical
section.  The invariant information was used before this ordinary
conversion, not inferred from it afterward.

For clarity, the norm does not need a Galois cover.  If
$\nu:\widehat X\to X$ has degree $d$ and
$s\in H^0(\widehat X,-kK_{\widehat X})$ is nonzero, \'etaleness gives
$-kK_{\widehat X}=\nu^*(-kK_X)$.
Locally the product of $s$ over the $d$ sheets is invariant under their
permutation, so it glues to a section of $-dkK_X$.
At a general point every factor is nonzero, so the norm is nonzero.
This is \cite[Lemma~2.2]{CompanionH}, here with all
its hypotheses visible.

\subsection{Recovering an invariant factor section}

There is also a useful converse transport after ordinary conversion.
It identifies precisely the analytic role of isometric monodromy.

\begin{proposition}\label{inv:liouville-recovery}
Let $d>0$ be an integer.  Suppose a compact complex manifold $X'$ is the quotient of
$\mathbb C^a\times W\times F$ by a group acting by product maps,
whose translation projection is a full lattice in $\mathbb C^a$.
Assume $W,F$ are connected compact complex manifolds, the canonical
bundles of $\mathbb C^a$ and $W$ have invariant nowhere-vanishing
frames, and monodromy on $H^0(F,-dK_F)$ preserves a Hermitian norm.
A nonzero section of $-dK_{X'}$ determines a nonzero
monodromy-invariant section of $-dK_F$.
\end{proposition}
\begin{proof}
Lift the section and trivialize its coefficient on the first two
factors using the invariant frames.  It becomes a family of sections
of $-dK_F$.  This family is independent of $W$: its coefficients are
holomorphic functions on a connected compact manifold and hence
constant.  More explicitly, choose finitely many evaluations at points
of $F$, in local frames, giving an injective linear map from the
finite-dimensional space $H^0(F,-dK_F)$ to a coordinate space.  These
evaluations also show that the family is a holomorphic map
\[
 v:\mathbb C^a\longrightarrow H^0(F,-dK_F).
\]
Deck invariance and the preserved Hermitian norm make $\|v(z)\|$
periodic under the full translation lattice.  It is bounded on a
fundamental parallelepiped and hence everywhere.  Liouville's theorem
makes every coordinate of $v$ constant.  Its nonzero constant value is
fixed by every projected deck transformation.  The same proof includes
$a=0$.
\end{proof}

\section{Extremal characters and invariant tensors}
\label{ext:section}

The tensor conversion proved above is ordinary: its input tensors need
not be invariant.  This permits another way to handle compact
monodromy.  We first construct invariant tensors on the compact factor,
then descend them before applying Corollary~\ref{conv:tensor-conversion}.
The character construction below uses the extreme terms of the full
equivariant Euler characteristic.  It is independent of the
weight-zero index calculation used in the preceding applications.

\begin{proposition}\label{ext:invariant-tensors}
Let $F$ be a smooth connected projective complex variety such that
$H^0(F,\Omega_F^i)=0$ for every $i>0$.  Suppose $L_F=-K_F$ has a
smooth semipositive Hermitian metric.  Let an algebraic torus $T$ act
effectively on $F$, and give $L_F$ its natural tangent-determinant
linearization.  Then there are an integer $p\geq0$ and unbounded
positive integers $k$ such that
\[
 H^0\bigl(F,(\Omega_F^1)^{\otimes p}\otimes kL_F\bigr)^T\ne0.
\]
The trivial torus and the point variety are included.
\end{proposition}

We use the left action on sections: a group element transports a
fiber value and changes its argument by the inverse action.  Thus a
tangent weight $a$ induces conormal weight $-a$.  This convention
fixes both the orbit-degree identity and the expansions below.  The
use of torus characters follows the equivariant viewpoint in
\cite{Muller2025}; the present argument produces tensors to be
converted on the total space, without a semiampleness assumption on
the compact factor.

\begin{proof}
Assume first that $T$ is nontrivial.  Choose a general lattice
cocharacter $u:\mathbb C^*\to T$ for which $F^u=F^T$.  A
$T$-linearized projective embedding shows that such choices exist:
one avoids the finitely many hyperplanes given by differences of its
weights.  The fixed locus has finitely many smooth connected
components $P$.  Write $w_P\in X^*(T)$ for the character of
$L_F|_P$, and put $\lambda_P=\langle w_P,u\rangle$.

\emph{The source has maximal anticanonical weight.}
The Bia\l ynicki-Birula decomposition~\cite{BB1973} has a unique
component $P_s$ whose attracting set as $t\to0$ is open.  The
attracting and repelling sets are affine-space bundles with ranks
given by the positive and negative normal weights; see also
\cite[Theorem~1.5 and Corollary~7.3]{JelisiejewSienkiewicz2019}.
Every normal weight at $P_s$ is positive on $u$.  The component is
proper because the action is nontrivial and effective.  The natural
anticanonical character is the sum of its normal tangent characters,
so
\begin{equation}\label{ext:source-positive}
 \lambda_{P_s}>0.
\end{equation}
For every other component there is a negative normal weight.  Its
repelling cell supplies a nonconstant orbit with limit in that
component as $t\to\infty$.  Extend the \emph{parametrized} orbit to
$a:\mathbb P^1\to F$ by properness.  For a linearized line $B$ the
fiber weights at its two endpoints satisfy
\begin{equation}\label{ext:orbit-degree}
 \operatorname{wt}_u(B|_{a(0)})-
 \operatorname{wt}_u(B|_{a(\infty)})
       =\deg a^*B.
\end{equation}
For example this follows by comparing equivariant frames at zero
and infinity on $\mathbb P^1$.  The parametrization, rather than
only the normalization of its image, includes any finite stabilizer
multiplicity.  For an ample linearized $B$ the difference is
positive, and for the nef line $L_F$ it is nonnegative.  Iterating
backward through fixed components terminates, because the ample
weights strictly increase in a finite set.  It can terminate only
at the source.  Therefore
\begin{equation}\label{ext:source-maximum}
 \lambda_{P_s}\geq\lambda_P\quad\hbox{for all }P.
\end{equation}

The attracting-cell limit map is a dominant rational map
$F\dashrightarrow P_s$.  Resolve it.  Pullback of holomorphic
forms along a dominant map is injective in characteristic zero,
and holomorphic forms are unchanged under smooth birational
modification.  Consequently $P_s$ has no positive-degree
holomorphic forms.  Since $P_s$ is connected and projective,
Hodge symmetry gives
\begin{equation}\label{ext:source-euler}
 \chi(P_s,\mathcal O_{P_s})=1.
\end{equation}

\emph{Localization detects the complete source character.}
Apply fixed-component localization to the virtual character
\[
 \chi_T(F,\ell L_F)=
       \sum_q(-1)^q[H^q(F,\ell L_F)]
\]
for positive integers $\ell$; we use the holomorphic index formula
of \cite{AtiyahSegal1968,AtiyahSinger1968}.  On a fixed component,
write a tangent normal block as $N_a e^a$, where $N_a$ is its
underlying rank-$r$ bundle with the weight removed.  Its contribution
is the factor $\lambda_{-1}(N_a^*e^{-a})^{-1}$.  Expand in the
direction of decreasing pairing with $u$.  If
$\langle a,u\rangle>0$, this factor is
\[
 \sum_{j\geq0}\operatorname{Sym}^j(N_a^*)e^{-ja};
\]
the leading term is $1$ and all others have negative pairing.
If $\langle a,u\rangle<0$, the same factor is
\[
 (-1)^r\det(N_a)e^{ra}
       \sum_{j\geq0}\operatorname{Sym}^j(N_a)e^{ja};
\]
even its leading term has strictly negative pairing.  These
identities follow by factoring the highest exterior term in the
second case.  Every normal character pairs nontrivially with $u$,
so each fixed exponent receives only finitely many contributions;
equivalently one works in the descending Laurent completion over
the character field of the sublattice pairing to zero.

By \eqref{ext:source-maximum}, every nonsource contribution is
strictly below $\ell\lambda_{P_s}$, including components with
$\lambda_P=\lambda_{P_s}$.  At the source the only term at that
pairing is the constant term of each positive normal factor.
Thus the entire part at the highest pairing is
\begin{equation}\label{ext:highest-character}
 e^{\ell w_{P_s}}\,
       \chi(P_s,\ell L_F|_{P_s}).
\end{equation}
This records the full $T$-character, not merely its restriction
to the chosen one-parameter subgroup.  Riemann--Roch makes its
coefficient a polynomial in $\ell$ whose value at zero is $1$
by \eqref{ext:source-euler}.  The coefficient is therefore
nonzero for all sufficiently large positive $\ell$.  Hence the
character $\ell w_{P_s}$ occurs in at least one actual cohomology
group $H^q(F,\ell L_F)$ for each such $\ell$.

\emph{Canceling source characters produces tensors.}
Only finitely many fixed components can occur as sources while
$u$ varies among general lattice directions.  Zero lies in the
convex hull of their characters.  Otherwise strict rational
separation, followed by a small general perturbation, gives a
lattice direction pairing negatively with every source character.
Its own source contradicts \eqref{ext:source-positive}.  Since
the characters are integral, there are nonnegative integers $b_s$,
not all zero, with
\begin{equation}\label{ext:cancel}
 \sum_s b_sw_{P_s}=0.
\end{equation}
For each source with $b_s>0$, choose a general direction having
that source.  Formula~\eqref{ext:highest-character} applies for
all sufficiently large $\ell$ in each of these finitely many
choices.  Passing to an infinite common subsequence fixes the
cohomological degrees $q_s$ in which the characters occur.

Average a K\"ahler form and the logarithmic weights of the given
line metric over the maximal compact torus $T_c$.  Smooth
semipositivity is preserved.  Hard Lefschetz with the smooth
coefficient $(\ell+1)L_F$~\cite[Theorem~0.1]{DPS2001} gives
$T_c$-equivariant surjections
\begin{equation}\label{ext:lefschetz}
 H^0\bigl(F,\Omega_F^{h-q_s}\otimes(\ell+1)L_F\bigr)
 \longrightarrow H^{q_s}(F,\ell L_F),\qquad h=\dim F.
\end{equation}
The target identification is equivariant because
$K_F+(\ell+1)L_F=\ell L_F$ with the natural linearizations.
Projecting a lift to its character space by weighted Haar averaging
therefore lifts the exact character $\ell w_{P_s}$.

Tensor the lifted nonzero sections with multiplicities $b_s$ and
use the alternating inclusions
$\Omega_F^j\hookrightarrow(\Omega_F^1)^{\otimes j}$ in
characteristic zero.  Their tensor product is nonzero at the generic
point, since $F$ is integral.  By \eqref{ext:cancel} its character
is zero.  Its tensor degree and line exponent are
\[
 p=\sum_s b_s(h-q_s),\qquad
 k=(\ell+1)\sum_s b_s.
\]
The first is fixed and the second tends to infinity.  Invariance
under $T_c$ is equivalent to $T$-invariance in the finite-dimensional
algebraic section representation.  This proves the nontrivial-torus
case.

If $T$ is trivial, Hodge symmetry gives
$\chi(F,\mathcal O_F)=1$.  Riemann--Roch for $\ell L_F$ is a
polynomial with this constant term, so some cohomology group is
nonzero for all sufficiently large $\ell$.  Fix its degree along
an infinite subsequence and apply \eqref{ext:lefschetz}.
For a point take $p=0$ and the constant section for every $k>0$.
\end{proof}

\subsection{Descent before ordinary conversion}

We finish this method by specifying the tensor descent.  Use the
finite-cover structure statement as in
Section~\ref{inv:isometric-descent}: on the universal cover of a
connected finite \'etale cover $\widehat X\to X$ there is an
isometric holomorphic product
\[
 \widetilde X=\mathbb C^a\times R\times F.
\]
Here $R$ is the compact Ricci-flat product, $F$ is projective with
no positive-degree holomorphic forms, and $-K_F$ is smoothly
semipositive.  The chosen deck subgroup acts by translations on
the flat factor, fixes the canonical frame on $R$, and acts on
$F$ through a connected compact torus of holomorphic isometries.
Point factors are allowed.  These are the geometric structure and
frame assertions, independent of the anticanonical nonvanishing
conclusion.

For completeness, the compact action has the algebraic realization
needed by Proposition~\ref{ext:invariant-tensors}.
Since $H^1(F,\mathcal O_F)=0$, it preserves a very ample line up
to isomorphism and hence acts projectively in its complete
linear-system embedding.  The representation is continuous, as
one sees by tracking a projective frame of points on the embedded
variety.  A connected compact abelian subgroup of the projective
linear group is conjugate into a diagonal torus.  Its Zariski
closure is an algebraic torus preserving $F$; quotient by its
ineffective kernel if necessary.  A diagonal projective torus
has a linear lift after fixing one diagonal coordinate, so an
ample linearization is available.  The action on $-K_F$ remains
the intrinsic tangent-determinant action.

Pull each invariant tensor $u_k$ from the proposition to
$\widetilde X$ and use the cotangent inclusion from the product
projection to $F$.  If $\omega_{\mathrm{flat}}$ and $\omega_R$
are the invariant canonical frames of the other factors, then
\[
 u_k\otimes(\omega_{\mathrm{flat}}\otimes\omega_R)^{-k}
 \in H^0\bigl(\widetilde X,
          (\Omega_{\widetilde X}^1)^{\otimes p}
                         \otimes(-kK_{\widetilde X})\bigr)
\]
is nonzero and deck-invariant.  It descends to $\widehat X$ and
is algebraic by projectivity.  Corollary~\ref{conv:tensor-conversion}
now gives a nonzero anticanonical plurisection on this actual
projective finite cover.  Finally the finite \'etale norm recalled
in Section~\ref{inv:isometric-descent} gives a nonzero positive
anticanonical multiple on $X$.  Thus the only invariant objects
needed for this method are the tensors before descent; the
conversion and its Fano-type auxiliary model are ordinary.

\bibliographystyle{amsplain}
\bibliography{references}
\end{document}
